% concatenated from shuffle_asymptotics_part3.tex
\documentclass{article}
\usepackage{amsmath,amsthm,amssymb,mathtools}
\usepackage{hyperref}
\hypersetup{
  hidelinks,
  pdfauthor={Alex Shvets},
  pdftitle={Universal Shuffle Asymptotics, Part III: Dominant-Block Quotient Geometry and Hybrid Gaussian--Compound-Poisson Limits in Finite-Alphabet Shuffle Privacy}
}
\usepackage{geometry}
\geometry{margin=1in}
\usepackage[most]{tcolorbox}
\tcbset{
  keybox/.style={
    enhanced,
    colback=gray!8,
    colframe=gray!50,
    boxrule=0.5pt,
    arc=2pt,
    left=6pt, right=6pt, top=6pt, bottom=6pt,
    fontupper=\small
  }
}

\newtheorem{theorem}{Theorem}[section]
\newtheorem{proposition}[theorem]{Proposition}
\newtheorem{lemma}[theorem]{Lemma}
\newtheorem{corollary}[theorem]{Corollary}
\newtheorem{definition}[theorem]{Definition}
\newtheorem{remark}[theorem]{Remark}

\newcommand{\TV}{\mathrm{TV}}
\newcommand{\Poi}{\mathrm{Poi}}
\newcommand{\Bin}{\mathrm{Bin}}
\newcommand{\Mult}{\mathrm{Mult}}
\newcommand{\Cat}{\mathrm{Cat}}
\newcommand{\LC}{\Delta_{\mathrm{LC}}}
\newcommand{\E}{\mathbb E}
\newcommand{\Pbb}{\mathbb P}
\newcommand{\R}{\mathbb R}
\newcommand{\N}{\mathbb N}
\newcommand{\1}{\mathbf 1}
\newcommand{\eps}{\varepsilon}
\newcommand{\weakto}{\Longrightarrow}
\newcommand{\privacycurve}{\mbox{privacy-curve}}
\DeclareMathOperator{\supp}{supp}
\DeclareMathOperator{\Var}{Var}
\DeclareMathOperator{\Cov}{Cov}

\title{Universal Shuffle Asymptotics, Part III:\\
Dominant-Block Quotient Geometry and Hybrid Gaussian--Compound-Poisson Limits\\
in Finite-Alphabet Shuffle Privacy}
\author{Alex Shvets\thanks{Email: \href{mailto:alt178332@gmail.com}{alt178332@gmail.com}. ORCID: \href{https://orcid.org/0009-0005-9802-379X}{0009-0005-9802-379X}.}}
\date{March 2026}

\begin{document}

\maketitle

\begin{abstract}

Part~I of this series~\cite{Shv26I} establishes a sharp Gaussian (LAN/GDP, in the sense of \cite{DRS22})
limit theory for neighboring shuffle experiments in the fixed full-support regime. Part~II~\cite{Shv26II}
identifies the first universality-breaking frontier: when the local randomizer becomes increasingly concentrated,
the Gaussian limit fails and the neighboring shuffle experiment enters critical Poisson, Skellam, and multivariate
compound-Poisson regimes.

The present paper completes the finite-alphabet weak-limit theory by identifying the
dominant-block quotient geometry that governs neighboring shuffle experiments. We treat
dominant blocks of arbitrary finite size, allow overlap between the dominant output sets
under the two neighboring hypotheses, and show that the limiting experiment decomposes
according to this geometry: projecting onto the sum of the dominant tangent spaces yields
a Gaussian factor, while quotienting by those same tangent spaces isolates a
compound-Poisson jump field in the rare block, recovering the critical Poisson, Skellam,
and multivariate compound-Poisson regimes of Part~II in the projected quotient block.
We also identify the regimes in which this quotient description determines the full
\privacycurve, as well as the obstruction that appears when projected jump limits alone
do not suffice.

Two further sections sharpen the rate picture and the boundary interface. First, we show that the
$O(n^{-1/2})$ rate for the full hybrid experiment is sharp in general, by an explicit binomial-plus-Bernoulli
calculation, and we identify a compatibility condition under which the sharper $O(n^{-1})$ rate is restored.
Second, we prove a boundary Berry--Esseen theorem showing that the critical Poisson-shift experiment is within
$O(c)$ in Le Cam distance of the Gaussian shift experiment as the critical scale parameter $c\downarrow 0$,
with a matching \privacycurve\ corollary. We also record a strong-boundary obstruction showing that, without an
additional structural hypothesis, projected jump limits need not determine the full \privacycurve. Together with
Parts~I--II, this yields a three-regime universality picture and a precise finite-alphabet L\'evy--Khintchine
layer for shuffle privacy.

\end{abstract}

\noindent\textbf{MSC 2020:} 62B15 (statistical experiments and information); 68P27 (privacy); 60F05 (central limit and other weak theorems); 60E07 (infinitely divisible distributions).\\[2pt]
\noindent\textbf{Keywords:} shuffle model, differential privacy, Le Cam distance, L\'evy--Khintchine formula, compound Poisson, multinomial normal approximation.

\tableofcontents

\section{Introduction}

Part~I~\cite{Shv26I} develops the Gaussian regime for neighboring shuffle experiments: for a fixed finite-output
local randomizer with full support bounded away from zero, the released histogram admits a complete
conditional-expectation linearization, the neighboring experiment is asymptotically Gaussian in the sense of
Le Cam, and the \privacycurve\ converges to the GDP curve~\cite{DRS22} with Berry--Esseen control.
Part~II~\cite{Shv26II} shows that this Gaussian picture breaks at the critical concentration threshold.
In the canonical one-dominant setting the neighboring experiment converges to a Poisson-shift limit; for
proportional compositions it converges to a Skellam-shift limit; and for general finite alphabets in the
sparse-error regime the released histogram converges, after centering, to a multivariate compound-Poisson law.

The shuffle model was initiated by \cite{EFM19,CSU19} and further developed in \cite{BBG19,FMT21,FMT23}.
For background on Le Cam theory and asymptotic statistics, see \cite{LeC86,vdV98}; for Poisson
approximation tools used repeatedly below, see \cite{BHJ92,Roos19}.

The remaining gap is structural rather than scalar. Part~II isolates the first hybrid weak limit in the
special two-dominant disjoint case, where the dominant block carries a Gaussian factor and the rare block
carries a compound-Poisson factor, but it does not cover overlapping dominant sets, dominant blocks of size
larger than two, or the general quotient geometry induced by several comparable dominant outputs. This paper
closes that gap at the level of weak limits and projected experiments. The central theorem identifies the
dominant tangent space
\[
M=M_0+M_1,
\]
projects the released histogram onto the Gaussian block $M$, quotients out the same tangent directions to the
rare block $M^\perp$, and proves a general finite-alphabet L\'evy--Khintchine limit
\[
(G,J)\quad\text{versus}\quad (G,J+\Delta)
\]
for the neighboring shuffled experiment. The overlap case does not require a different limiting object: the same
hybrid limit persists, but the deterministic quotient shift $\Delta$ may collapse to zero.

The main results are Theorem~\ref{thm:general-k-dominant-LK}, which gives weak convergence of the full hybrid
statistic, $O(n^{-1})$ total-variation and Le Cam convergence for the projected jump experiment, and
\privacycurve\ convergence for the full experiment in the interior regime, in the weak boundary regime, and in
the regular strong-boundary cases in which no unsmoothed finite minority tangent block remains. A short
counterexample shows that no fully general strong-boundary \privacycurve\ theorem is possible: a bounded
minority block can retain dominant internal geometry that vanishes weakly but remains visible to exact histogram
tests. The theorem specializes back to the single-dominant multivariate compound-Poisson limit of Part~II and
to the two-dominant disjoint hybrid limit of Part~II, Appendix~B.

At the level of general theory, the abstract existence of infinitely divisible limit experiments is classical.
Strasser's scale-invariance theorem~\cite{Str85} and the monograph of Janssen, Milbrodt, and Strasser~\cite{JMS85}
develop a general L\'evy--Khintchine theory for statistical experiments with independent increments. What is new
here is the explicit finite-alphabet shuffle geometry: the dominant tangent space, the quotient by the dominant
hypergraph, the compound-Poisson jump field attached to rare outputs, and the accompanying \privacycurve\
statements. On the Gaussian side, Carter's deficiency bound between multinomial and multivariate normal
experiments~\cite{Car02} is a central precursor, but it treats the dense regime in which all relevant cell
probabilities stay bounded away from zero; our setting contains genuinely rare $O(n^{-1})$ cells and therefore
falls outside that framework. At the level of shuffle privacy, Takagi and Liew~\cite{TL26} obtain asymptotic
privacy amplification bounds beyond pure local differential privacy, while Daskalakis, Kamath, and
Tzamos~\cite{DKT15} give a structural Gaussian-plus-sparse decomposition for Poisson multinomial distributions.
The present paper is complementary to both lines: it works at the level of neighboring binary experiments,
explicit Le Cam limits, and sharp \privacycurve\ asymptotics for shuffled histograms.

\section{Model and notation}

\subsection{Shuffle mechanism, transcript laws, and privacy curves}
\label{sec:model-privacy-curves}

Fix a population size $n\ge 1$. Each user $i\in\{1,\dots,n\}$ holds a private datum $x_i\in\{0,1\}$. A
(possibly $n$-dependent) local randomizer is a Markov kernel
\[
W^{(n)}:\{0,1\}\to \Delta(\mathcal Y),\qquad b\mapsto W_b^{(n)},
\]
where $\mathcal Y$ is a finite output alphabet and $\Delta(\mathcal Y)$ denotes the simplex of probability
measures on $\mathcal Y$. Given an input dataset $x^n=(x_1,\dots,x_n)$, users apply $W^{(n)}$ independently to
produce messages $Y_i\sim W^{(n)}_{x_i}$. The shuffle mechanism outputs the multiset of messages, equivalently
its histogram
\[
N(y):=\sum_{i=1}^n \1\{Y_i=y\}\in \mathbb Z^{\mathcal Y},\qquad \sum_{y\in\mathcal Y}N(y)=n.
\]
For $k\in\{0,\dots,n\}$, let $T_{n,k}$ denote the law of the released histogram when the dataset has exactly $k$
ones and $n-k$ zeros.

The neighboring shuffle experiment associated with $k$ versus $k+1$ ones is the binary experiment
\[
\mathcal E_{n,k}:=(P_{n,k},Q_{n,k}),\qquad P_{n,k}:=T_{n,k},\qquad Q_{n,k}:=T_{n,k+1}.
\]

For a binary experiment $(P,Q)$ on any measurable space, the one-sided privacy curve is
\[
\delta_{Q\|P}(\varepsilon)
:=
\sup_A\bigl\{Q(A)-e^\varepsilon P(A)\bigr\},
\qquad \varepsilon\ge 0,
\]
where the supremum ranges over all measurable sets $A$.

\begin{lemma}[Contraction of total variation under measurable maps]
\label{lem:tv-contraction-main}
Let $\mu,\nu$ be probability measures on $(\mathcal X,\mathcal F)$ and let $f:\mathcal X\to\mathcal Z$ be measurable.
Then
\[
\TV(\mu\circ f^{-1},\nu\circ f^{-1})\le \TV(\mu,\nu).
\]
\end{lemma}

\begin{proof}
For any measurable $B\subseteq\mathcal Z$,
\[
(\mu\circ f^{-1})(B)-(\nu\circ f^{-1})(B)=\mu(f^{-1}(B))-\nu(f^{-1}(B)).
\]
Taking absolute values and then the supremum over $B$ yields the claim.
\end{proof}

\begin{lemma}[Tensorization bound for product measures]
\label{lem:tv-tensorization-main}
If $\mu_i,\nu_i$ are probability measures on measurable spaces $(\mathcal X_i,\mathcal F_i)$, then
\[
\TV(\mu_1\otimes\mu_2,\nu_1\otimes\nu_2)
\le
\TV(\mu_1,\nu_1)+\TV(\mu_2,\nu_2).
\]
More generally, for finitely many factors,
\[
\TV\Bigl(\bigotimes_{j=1}^m \mu_j,\bigotimes_{j=1}^m \nu_j\Bigr)
\le \sum_{j=1}^m \TV(\mu_j,\nu_j).
\]
\end{lemma}

\begin{proof}
Choose independent couplings $(X_i,Y_i)$ of $(\mu_i,\nu_i)$ such that
\[
\Pbb(X_i\neq Y_i)=\TV(\mu_i,\nu_i).
\]
Then $(X_1,X_2)$ and $(Y_1,Y_2)$ form a coupling of $\mu_1\otimes\mu_2$ and $\nu_1\otimes\nu_2$, hence
\[
\TV(\mu_1\otimes\mu_2,\nu_1\otimes\nu_2)
\le
\Pbb((X_1,X_2)\neq (Y_1,Y_2))
\le \Pbb(X_1\neq Y_1)+\Pbb(X_2\neq Y_2).
\]
The general case follows by induction.
\end{proof}

\begin{lemma}[Le Cam distance on the same space]
\label{lem:lecam-tv-bound-main}
Let $(P,Q)$ and $(P',Q')$ be binary experiments on the same measurable space. Then
\[
\LC\bigl((P,Q),(P',Q')\bigr)
\le
\max\{\TV(P,P'),\TV(Q,Q')\}.
\]
\end{lemma}

\begin{proof}
Choose the identity Markov kernel in both deficiencies. Each deficiency is then bounded by the corresponding
total-variation distance, and taking the maximum gives the claim.
\end{proof}

\begin{lemma}[Privacy-curve stability under total-variation perturbations]
\label{lem:privacy-stability-main}
Let $(P,Q)$ and $(P',Q')$ be binary experiments on the same measurable space. Then for every $\varepsilon\ge 0$,
\[
\bigl|\delta_{Q\|P}(\varepsilon)-\delta_{Q'\|P'}(\varepsilon)\bigr|
\le
\TV(Q,Q')+e^\varepsilon\TV(P,P').
\]
\end{lemma}

\begin{proof}
For every measurable set $A$,
\[
Q(A)-e^\varepsilon P(A)
\le
Q'(A)+\TV(Q,Q')-e^\varepsilon P'(A)+e^\varepsilon\TV(P,P').
\]
Taking the supremum over $A$ yields
\[
\delta_{Q\|P}(\varepsilon)
\le
\delta_{Q'\|P'}(\varepsilon)+\TV(Q,Q')+e^\varepsilon\TV(P,P').
\]
Exchanging $(P,Q)$ and $(P',Q')$ gives the reverse inequality.
\end{proof}


\begin{lemma}[Data processing for privacy curves]
\label{lem:privacy-dpi-main}
Let $(P,Q)$ be a binary experiment on $\mathcal X$, let $K$ be a Markov kernel from $\mathcal X$ to
$\mathcal Z$, and write
\[
P':=PK,
\qquad
Q':=QK.
\]
Then for every $\varepsilon\ge 0$,
\[
\delta_{Q'\|P'}(\varepsilon)\le \delta_{Q\|P}(\varepsilon).
\]
\end{lemma}

\begin{proof}
Set $\mu:=Q-e^\varepsilon P$, viewed as a finite signed measure on $\mathcal X$. For any measurable
$B\subseteq\mathcal Z$, the function $f_B(x):=K(x,B)$ is measurable and takes values in $[0,1]$, so
\[
Q'(B)-e^\varepsilon P'(B)=\int f_B\,d\mu
\le \sup_{0\le f\le 1}\int f\,d\mu.
\]
By the Hahn decomposition theorem, the last supremum equals $\mu^+(\mathcal X)=\sup_A \mu(A)$, where the
supremum on the right ranges over measurable $A\subseteq\mathcal X$. Thus
\[
Q'(B)-e^\varepsilon P'(B)\le \sup_A\{Q(A)-e^\varepsilon P(A)\}
=\delta_{Q\|P}(\varepsilon).
\]
Taking the supremum over $B$ proves the claim.
\end{proof}

\begin{lemma}[Blackwell-equivalent experiments have identical privacy curves]
\label{lem:privacy-blackwell-equivalence}
Let $(P,Q)$ be a binary experiment on $\mathcal X$ and let $(P',Q')$ be a binary experiment on $\mathcal X'$.
Suppose there exist Markov kernels $K:\mathcal X\to\mathcal X'$ and $L:\mathcal X'\to\mathcal X$ such that
\[
PK=P',
\qquad
QK=Q',
\qquad
P'L=P,
\qquad
Q'L=Q.
\]
Then for every $\varepsilon\ge 0$,
\[
\delta_{Q\|P}(\varepsilon)=\delta_{Q'\|P'}(\varepsilon).
\]
\end{lemma}

\begin{proof}
Applying Lemma~\ref{lem:privacy-dpi-main} to $K$ gives
\[
\delta_{Q'\|P'}(\varepsilon)\le \delta_{Q\|P}(\varepsilon),
\]
and applying it to $L$ gives the reverse inequality.
\end{proof}

\begin{lemma}[Common independent factor does not affect privacy curves]
\label{lem:common-independent-factor-main}
Let $(P,Q)$ be a binary experiment on $\mathcal X$, and let $R$ be any probability measure on $\mathcal Z$.
Then for every $\varepsilon\ge 0$,
\[
\delta_{Q\otimes R\|P\otimes R}(\varepsilon)=\delta_{Q\|P}(\varepsilon).
\]
\end{lemma}

\begin{proof}
For any measurable $A\subseteq \mathcal X\times\mathcal Z$, define $A_z:=\{x\in\mathcal X:(x,z)\in A\}$. Then
\[
(Q\otimes R)(A)-e^\varepsilon(P\otimes R)(A)
=
\int \bigl(Q(A_z)-e^\varepsilon P(A_z)\bigr)\,R(dz)
\le \delta_{Q\|P}(\varepsilon).
\]
Taking the supremum over $A$ gives
\[
\delta_{Q\otimes R\|P\otimes R}(\varepsilon)\le \delta_{Q\|P}(\varepsilon).
\]
For the reverse bound, fix a measurable $B\subseteq\mathcal X$ and take $A=B\times\mathcal Z$. Then
\[
(Q\otimes R)(A)-e^\varepsilon(P\otimes R)(A)=Q(B)-e^\varepsilon P(B).
\]
Taking the supremum over $B$ proves equality.
\end{proof}

\subsection{General finite-dominant regime and quotient geometry}

\begin{definition}[General finite-dominant sparse-error regime]
\label{def:general-k-dominant}
Fix a finite output alphabet $\mathcal Y$ and write $(e_y)_{y\in\mathcal Y}$ for the standard basis of
$\mathbb R^{\mathcal Y}$. For each $n\ge 1$ and $b\in\{0,1\}$, let
\[
W_b^{(n)} \in \Delta(\mathcal Y)
\]
be the local randomizer law under input $b$.
For $k\in\{0,\dots,n\}$, let $T_{n,k}$ denote the shuffled transcript law for a dataset with
$n-k$ zeros and $k$ ones, and let
\[
N_{n,k}\in \mathbb Z^{\mathcal Y}
\]
be the released histogram under $T_{n,k}$.

We say that $W^{(n)}=(W_0^{(n)},W_1^{(n)})$ is in the
\emph{general finite-dominant sparse-error critical regime} if, for each $b\in\{0,1\}$, there exist

\begin{itemize}
\item a nonempty dominant set $D_b\subseteq \mathcal Y$,
\item a dominant law $p_b\in \Delta(D_b)$ with $p_b(y)>0$ for every $y\in D_b$,
\item rare intensities $\alpha_b(y)\in[0,\infty)$ for $y\in \mathcal Y\setminus D_b$,
\end{itemize}

such that
\begin{equation}
W_b^{(n)}(y)=p_b(y)+O(n^{-1})
\qquad (y\in D_b),
\label{eq:gen-dom-dominant}
\end{equation}
and
\begin{equation}
n\,W_b^{(n)}(y)\longrightarrow \alpha_b(y)
\qquad (y\in \mathcal Y\setminus D_b).
\label{eq:gen-dom-rare}
\end{equation}
Write
\[
\mu_b:=\sum_{y\in D_b} p_b(y)e_y\in \mathbb R^{\mathcal Y}.
\]
Finally, let $k_n\in\{0,\dots,n-1\}$ be any composition sequence with
\[
\pi_n:=\frac{k_n}{n}\longrightarrow \pi\in[0,1].
\]
\end{definition}

\begin{remark}
The positivity condition $p_b(y)>0$ ensures that $D_b$ contains only genuinely dominant outputs; outputs with vanishing limiting probability should be classified as rare.
\end{remark}

\begin{definition}[Dominant tangent space and hypergraph quotient]
\label{def:dominant-tangent-quotient}
For each $b\in\{0,1\}$ define the dominant tangent subspace
\[
M_b
:=\mathrm{span}\{e_y-\mu_b:\ y\in D_b\}
=
\Bigl\{x\in\mathbb R^{\mathcal Y}:\ \mathrm{supp}(x)\subseteq D_b,\ \sum_{y\in D_b}x(y)=0\Bigr\},
\]
and set
\[
M:=M_0+M_1\subseteq \mathbb R^{\mathcal Y}.
\]
Let $\Pi_G$ be the orthogonal projection onto $M$, and let
\[
\Pi_J:=I-\Pi_G
\]
be the orthogonal projection onto $M^\perp$.
Thus $M$ is the Gaussian dominant-tangent space and $M^\perp$ is the rare quotient space.

Let $\mathsf H$ be the hypergraph with vertex set $D_0\cup D_1$ and hyperedges $D_0,D_1$.
Let $\mathcal C$ denote the set of connected components of $\mathsf H$.
Then
\[
M^\perp\cap \mathbb R^{D_0\cup D_1}
=
\{x\in \mathbb R^{D_0\cup D_1}:\ x \text{ is constant on each } C\in\mathcal C\},
\]
equivalently
\[
\Pi_J e_y=\Pi_J e_{y'}
\quad\Longleftrightarrow\quad
y,y' \text{ belong to the same component } C\in\mathcal C.
\]
(see Lemma~\ref{lem:new-geometry} in Appendix~\ref{app:technical-lemmas} for the proof).
For each component $C\in\mathcal C$, write
\[
m_C:=\Pi_J e_y \qquad (y\in C),
\]
which is well-defined by the preceding display.
\end{definition}

\begin{remark}
We recall that the privacy curve $\delta_{Q\|P}(\varepsilon)$ was defined in
Section~\ref{sec:model-privacy-curves}.
\end{remark}

\section{The general L\'evy--Khintchine limit theorem}

\begin{theorem}[General $k$-dominant L\'evy--Khintchine limit for neighboring shuffle experiments]
\label{thm:general-k-dominant-LK}
Assume the setting of Definition~\ref{def:general-k-dominant}; the privacy curve is as defined in
Section~\ref{sec:model-privacy-curves}.
Define the centered histogram
\[
\widehat H_n
:=
N_{n,k_n}-(n-k_n)\mu_0-k_n\mu_1\in\mathbb R^{\mathcal Y},
\]
and the hybrid normalized statistic
\[
S_n
:=
\Bigl(n^{-1/2}\Pi_G\widehat H_n,\ \Pi_J\widehat H_n\Bigr)
\in M\times M^\perp.
\]
Under the neighboring alternative $T_{n,k_n+1}$ we use the \emph{same} centering
$(n-k_n)\mu_0+k_n\mu_1$.

For $b\in\{0,1\}$ define the dominant covariance operator
\[
\Gamma_b
:=
\sum_{y\in D_b} p_b(y)\,(e_y-\mu_b)(e_y-\mu_b)^\top
\qquad\text{on }\mathbb R^{\mathcal Y},
\]
and set
\[
\Sigma:=(1-\pi)\Gamma_0+\pi\Gamma_1
\qquad\text{(viewed as an operator on }M\text{)}.
\]
For each $b\in\{0,1\}$ and each $y\in\mathcal Y\setminus D_b$, define the jump vector
\[
j_{b,y}:=\Pi_J(e_y-\mu_b)\in M^\perp.
\]
Let $\nu$ be the finite measure on $M^\perp\setminus\{0\}$ given by
\[
\nu
:=
\sum_{\substack{y\notin D_0\\ j_{0,y}\neq 0}}
(1-\pi)\alpha_0(y)\,\delta_{j_{0,y}}
+
\sum_{\substack{y\notin D_1\\ j_{1,y}\neq 0}}
\pi\alpha_1(y)\,\delta_{j_{1,y}}.
\]
Let $G\sim N(0,\Sigma)$, let $J$ be an independent compound-Poisson random vector with L\'evy measure
$\nu$, equivalently
\[
J\overset d=
\sum_{y\notin D_0} U_y\,j_{0,y}
+
\sum_{y\notin D_1} V_y\,j_{1,y},
\]
where all coordinates are independent and
\[
U_y\sim \mathrm{Poi}((1-\pi)\alpha_0(y)),
\qquad
V_y\sim \mathrm{Poi}(\pi\alpha_1(y)).
\]
Set
\[
\Delta:=\Pi_J(\mu_1-\mu_0)\in M^\perp.
\]
Let $P_n$ be the law of $S_n$ under $T_{n,k_n}$, and let $Q_n$ be the law of the same statistic $S_n$
under $T_{n,k_n+1}$.

Then the following hold.

\begin{enumerate}
\item[(i)] \textbf{Weak L\'evy--Khintchine convergence of the full hybrid statistic.}
There exist probability laws
\[
P_\infty:=\mathcal L(G,J),
\qquad
Q_\infty:=\mathcal L(G,J+\Delta),
\]
such that
\[
P_n \Longrightarrow P_\infty,
\qquad
Q_n \Longrightarrow Q_\infty.
\]
Equivalently, for every $u\in M$ and $v\in M^\perp$,
\[
\mathbb E\exp\!\bigl(i\langle u,G\rangle+i\langle v,J\rangle\bigr)
=
\exp\!\left(
-\frac12\langle u,\Sigma u\rangle
+
\int_{M^\perp\setminus\{0\}}
\bigl(e^{i\langle v,z\rangle}-1\bigr)\,\nu(dz)
\right),
\]
whereas the limiting characteristic function under the neighboring alternative is multiplied by
$e^{i\langle v,\Delta\rangle}$.

In particular, under the neighboring calibration considered here, the Gaussian factor is asymptotically
\emph{common} to both hypotheses: no $M$-shift appears in the limit, and the entire neighboring shift
is carried by the quotient coordinate $\Delta\in M^\perp$.

Moreover, if
\[
\bar p_b(C):=\sum_{y\in D_b\cap C} p_b(y)
\qquad (C\in\mathcal C,\ b\in\{0,1\}),
\]
then
\[
\Delta
=
\sum_{C\in\mathcal C}\bigl(\bar p_1(C)-\bar p_0(C)\bigr)\,m_C.
\]
Indeed, for each $b\in\{0,1\}$,
\[
\Pi_J\mu_b
=\Pi_J\Bigl(\sum_{y\in D_b} p_b(y)e_y\Bigr)
=\sum_{y\in D_b} p_b(y)\Pi_J e_y
=\sum_{C\in\mathcal C}\Bigl(\sum_{y\in D_b\cap C} p_b(y)\Bigr)m_C
=\sum_{C\in\mathcal C} \bar p_b(C)\,m_C,
\]
so subtracting the two identities gives the displayed formula for $\Delta$.
Hence the overlap case is not a different limit structure: it is the same theorem with a possibly
degenerate shift. In particular, if $D_0\cap D_1\neq\varnothing$, then $\mathsf H$ has a single
component, so $\Delta=0$ and therefore
\[
P_\infty=Q_\infty.
\]

\item[(ii)] \textbf{Projected total-variation and Le Cam convergence on the quotient block.}
Define the projected statistic
\[
S_n^J:=\Pi_J\widehat H_n\in M^\perp,
\]
let $P_n^J$ and $Q_n^J$ denote its laws under $T_{n,k_n}$ and $T_{n,k_n+1}$, and set
\[
P_\infty^J:=\mathcal L(J),
\qquad
Q_\infty^J:=\mathcal L(J+\Delta).
\]
Assume in addition that
\begin{equation}
\max_{b\in\{0,1\}}
\left(
\sum_{y\in D_b}\bigl|W_b^{(n)}(y)-p_b(y)\bigr|
+
\sum_{y\notin D_b}\bigl|nW_b^{(n)}(y)-\alpha_b(y)\bigr|
\right)
+
|\pi_n-\pi|
\le \frac{C_0}{n}
\label{eq:quantitative-critical-regime}
\end{equation}
for all sufficiently large $n$ and some finite constant $C_0$.
Then there exists a finite constant $C=C(C_0,p_0,p_1,\alpha_0,\alpha_1,\pi,\mathcal Y)$ such that
\[
\mathrm{TV}(P_n^J,P_\infty^J)
+
\mathrm{TV}(Q_n^J,Q_\infty^J)
\le \frac{C}{n}
\]
for all sufficiently large $n$.
Consequently,
\[
\Delta_{\mathrm{LC}}\!\bigl((P_n^J,Q_n^J),(P_\infty^J,Q_\infty^J)\bigr)
\le
\max\bigl\{\mathrm{TV}(P_n^J,P_\infty^J),\,\mathrm{TV}(Q_n^J,Q_\infty^J)\bigr\}
\le \frac{C}{n},
\]
where $\Delta_{\mathrm{LC}}$ denotes Le Cam distance.

\item[(iii)] \textbf{\privacycurve\ convergence for the full hybrid experiment in the regular regimes.}
For every fixed $\varepsilon\ge 0$,
\[
\delta_{Q_\infty\|P_\infty}(\varepsilon)
=
\delta_{\mathcal L(J+\Delta)\,\|\,\mathcal L(J)}(\varepsilon),
\]
because $G$ is common and independent under both limiting hypotheses.

Moreover, the convergence
\[
\delta_{Q_n\|P_n}(\varepsilon)\longrightarrow \delta_{Q_\infty\|P_\infty}(\varepsilon)
\]
holds in each of the following cases:
\begin{enumerate}
\item[(a)] \emph{Interior regime:} assume in addition \eqref{eq:quantitative-critical-regime} and $\pi\in(0,1)$.
\item[(b)] \emph{Weak boundary regime:} $\pi=0$ with $k_n\to\infty$ and $k_n/n\to 0$, or symmetrically
$\pi=1$ with $n-k_n\to\infty$ and $(n-k_n)/n\to 0$.
\item[(c)] \emph{Strong boundary regularity:} assume in addition \eqref{eq:quantitative-critical-regime},
$\pi=0$, $k_n=O(1)$, and either $D_1\subseteq D_0$ or
$|D_1|=1$; or symmetrically $\pi=1$, $n-k_n=O(1)$, and either $D_0\subseteq D_1$ or $|D_0|=1$.
\end{enumerate}

If, in addition, \eqref{eq:quantitative-critical-regime} holds and $\pi\in(0,1)$, then there exists a finite constant
$C_{\mathrm{int}}=C_{\mathrm{int}}(p_0,p_1,\alpha_0,\alpha_1,\pi,\mathcal Y)$ such that
\[
\bigl|
\delta_{Q_n\|P_n}(\varepsilon)-\delta_{Q_\infty\|P_\infty}(\varepsilon)
\bigr|
\le
\frac{C_{\mathrm{int}}(1+e^\varepsilon)}{\sqrt n}
\]
for all sufficiently large $n$.

Without the additional strong-boundary regularity in \textup{(c)}, the projected jump limit need not govern the
full \privacycurve; see Proposition~\ref{prop:strong-boundary-obstruction}.
\end{enumerate}
\end{theorem}

\paragraph{Specializations.}
\begin{enumerate}
\item[\textup{(a)}] \textbf{Single-dominant regime.}
If $|D_0|=|D_1|=1$, say $D_0=\{y_0\}$ and $D_1=\{y_1\}$, then
\[
M_0=M_1=\{0\},\qquad M=\{0\},\qquad \Pi_J=I,\qquad G\equiv 0.
\]
Moreover,
\[
j_{0,y}=e_y-e_{y_0}\quad (y\neq y_0),
\qquad
j_{1,y}=e_y-e_{y_1}\quad (y\neq y_1),
\qquad
\Delta=e_{y_1}-e_{y_0}.
\]
Hence part~\textup{(i)} reduces exactly to the multivariate compound-Poisson / Poisson-shift limit
of Theorem~5.8 of Part~II, and part~\textup{(ii)} reduces to its $O(n^{-1})$ total-variation / Le Cam
convergence statement.

\item[\textup{(b)}] \textbf{Two-dominant disjoint regime.}
If $|D_b|=2$ for each $b$ and $D_0\cap D_1=\varnothing$, write
\[
D_b=\{y_{ba},y_{bb}\},
\qquad
p_b(y_{ba})=p_b,\quad p_b(y_{bb})=1-p_b,
\qquad
g_b:=e_{y_{ba}}-e_{y_{bb}}.
\]
Then
\[
M=\mathrm{span}\{g_0,g_1\},
\qquad
\Gamma_b=p_b(1-p_b)\,g_bg_b^\top,
\qquad
\Sigma=(1-\pi)p_0(1-p_0)\,g_0g_0^\top+\pi p_1(1-p_1)\,g_1g_1^\top,
\]
and the theorem reduces to Proposition~5.4 of Part~II.
Its projected statement \textup{(ii)} is the natural $k$-dominant extension of the projected Le Cam
result of Section~5.3 of Part~II. In the interior and weak-boundary regimes, part~\textup{(iii)}
reduces to the corresponding full hybrid \privacycurve\ convergence statement of Appendix~B of Part~II.
At a strong boundary with a bounded minority group, however, an additional regularity condition is needed;
see part~\textup{(iii)} and Proposition~\ref{prop:strong-boundary-obstruction}.

\item[\textup{(c)}] \textbf{Two-dominant overlap regime.}
If $|D_b|=2$ for each $b$ and $D_0\cap D_1\neq\varnothing$, then $\mathsf H$ has a single connected
component, hence
\[
\Delta=0.
\]
Therefore the L\'evy--Khintchine weak limit of part~\textup{(i)} is common under both hypotheses:
\[
P_\infty=Q_\infty=\mathcal L(G,J).
\]
Thus overlap does \emph{not} require a different limiting object; it collapses the deterministic
jump-shift in the quotient block. At a strong boundary, overlap alone is not sufficient for
\privacycurve\ convergence unless the bounded minority dominant block is absorbed by the majority block;
see Proposition~\ref{prop:strong-boundary-obstruction}.

\item[\textup{(d)}] \textbf{Mixed-size dominant sets.}
If $|D_0|=1$ and $|D_1|\ge 2$ with $D_0\cap D_1=\varnothing$, then $M_0=\{0\}$,
$M=M_1$, and the Gaussian factor operates only on the tangent space of the larger
dominant set. The projected jump experiment retains the full Poisson-shift structure
from the single-dominant side, but at the strong boundary $\pi=0$ with $k_n=O(1)$ the full
\privacycurve\ need not be governed by that projected limit; this mixed-size case lies outside
the regular strong-boundary regime of part~\textup{(iii)}.
\end{enumerate}

\section{Proof of part (i): weak L\'evy--Khintchine convergence}
\label{sec:proof-i}

Throughout this section write
\[
m_{0,n}:=n-k_n,\qquad m_{1,n}:=k_n,
\]
and define
\[
G_n:=n^{-1/2}\Pi_G\widehat H_n,\qquad Z_n:=\Pi_J\widehat H_n,
\qquad S_n=(G_n,Z_n).
\]

\begin{proof}[Proof of Theorem~\ref{thm:general-k-dominant-LK}(i)]
\textbf{Step 1: establish the quotient geometry.}
For each $b\in\{0,1\}$ and $y\in D_b$, the vector $e_y-\mu_b$ belongs to $M_b\subseteq M$, hence
\[
\Pi_J(e_y-\mu_b)=0,
\qquad \Pi_J e_y=\Pi_J\mu_b=:m_b.
\]
Thus all dominant outputs inside the same dominant set collapse to a single quotient atom. This is
Lemma~\ref{lem:new-geometry} proved in Appendix~\ref{app:technical-lemmas}. It is the exact replacement for the
$m_0\neq m_1$ argument in Part~II, Corollary~5.9: the point is not that the dominant quotient atoms must be
distinct, but rather that the entire dominant block is annihilated by $\Pi_J$. If $D_0\cap D_1\neq\varnothing$, then
Lemma~\ref{lem:new-geometry} shows that $m_0=m_1$ and therefore
\[
\Delta=\Pi_J(\mu_1-\mu_0)=0.
\]

\medskip
\textbf{Step 2: rewrite the statistic as a triangular-array sum.}
Under $T_{n,k_n}$ the released histogram is
\[
N_{n,k_n}=\sum_{i=1}^{m_{0,n}} e_{Y^{(0)}_{i,n}}+\sum_{j=1}^{m_{1,n}} e_{Y^{(1)}_{j,n}},
\]
where the two groups are independent and $Y^{(b)}_{\ell,n}\sim W_b^{(n)}$. Set
\[
X^{(b)}_{\ell,n}:=e_{Y^{(b)}_{\ell,n}}-\mu_b,
\qquad
\xi^{(b)}_{\ell,n}:=\bigl(n^{-1/2}\Pi_G X^{(b)}_{\ell,n},\ \Pi_J X^{(b)}_{\ell,n}\bigr).
\]
Then
\[
S_n=\sum_{i=1}^{m_{0,n}}\xi^{(0)}_{i,n}+\sum_{j=1}^{m_{1,n}}\xi^{(1)}_{j,n}.
\]
Hence $S_n$ is a sum of independent increments in the triangular-array form required for a characteristic-
function argument.

\medskip
\textbf{Step 3: expand the one-user characteristic factor.}
Fix $u\in M$ and $v\in M^\perp$, and define
\[
\phi_{b,n}(u,v)
:=
\E\exp\!\left(
 i\bigl\langle u,n^{-1/2}\Pi_G X^{(b)}_{1,n}\bigr\rangle
 +i\bigl\langle v,\Pi_J X^{(b)}_{1,n}\bigr\rangle
\right).
\]
By Lemma~\ref{lem:new-one-user-cf} of Appendix~\ref{app:technical-lemmas}, uniformly on compact subsets of
$M\times M^\perp$,
\[
\log\phi_{b,n}(u,v)
=
-\frac{1}{2n}\langle u,\Gamma_b u\rangle
+
\frac{1}{n}\sum_{y\notin D_b}\alpha_b(y)\bigl(e^{i\langle v,j_{b,y}\rangle}-1\bigr)
+o(n^{-1}).
\]
Compared with Part~II, Proposition~5.4, the only change is that the dominant block is now genuinely
multivariate. The first-order term still vanishes because
\[
\sum_{y\in D_b} p_b(y)(e_y-\mu_b)=0,
\]
while the quadratic term becomes the covariance form $\langle u,\Gamma_bu\rangle$ rather than a scalar
binomial variance. The rare outputs still appear with probability $O(n^{-1})$, so they contribute only the
compound-Poisson exponent in the quotient coordinate. No mixed $u$--$v$ term survives, because the $u$-dependence
of each rare term is only $O(n^{-1/2})$ and is multiplied by a probability $O(n^{-1})$.

\medskip
\textbf{Step 4: sum logarithms and identify the limit under $T_{n,k_n}$.}
By independence across users,
\[
\log \E e^{i\langle u,G_n\rangle+i\langle v,Z_n\rangle}
=
 m_{0,n}\log\phi_{0,n}(u,v)+m_{1,n}\log\phi_{1,n}(u,v).
\]
Substituting the expansion from Step~3 and using $m_{0,n}/n\to 1-\pi$, $m_{1,n}/n\to \pi$, we obtain
\begin{align*}
\log \E e^{i\langle u,G_n\rangle+i\langle v,Z_n\rangle}
&\to
-\frac12\bigl((1-\pi)\langle u,\Gamma_0u\rangle+\pi\langle u,\Gamma_1u\rangle\bigr) \\
&\qquad +\sum_{y\notin D_0}(1-\pi)\alpha_0(y)\bigl(e^{i\langle v,j_{0,y}\rangle}-1\bigr)
+\sum_{y\notin D_1}\pi\alpha_1(y)\bigl(e^{i\langle v,j_{1,y}\rangle}-1\bigr).
\end{align*}
The right-hand side is exactly
\[
-\frac12\langle u,\Sigma u\rangle
+
\int_{M^\perp\setminus\{0\}} \bigl(e^{i\langle v,z\rangle}-1\bigr)\,\nu(dz),
\]
namely the characteristic exponent of an independent pair $(G,J)$ with $G\sim N(0,\Sigma)$ and $J$ compound-
Poisson with L\'evy measure $\nu$. L\'evy's continuity theorem therefore yields
\[
P_n\weakto \mathcal L(G,J)=P_\infty.
\]

\medskip
\textbf{Step 5: treat the neighboring alternative under the same centering.}
Under $T_{n,k_n+1}$, let $\widetilde S_n$ denote the hybrid statistic built from the actual group sizes
$(m_{0,n}-1,m_{1,n}+1)$ and centered by $(m_{0,n}-1)\mu_0+(m_{1,n}+1)\mu_1$. Then the statistic $S_n$ from the
statement satisfies
\[
S_n=\widetilde S_n+\bigl(n^{-1/2}\Pi_G(\mu_1-\mu_0),\ \Pi_J(\mu_1-\mu_0)\bigr)
\]
under $T_{n,k_n+1}$. The same calculation as in Step~4 applies to $\widetilde S_n$, because replacing
$m_{0,n},m_{1,n}$ by $m_{0,n}-1,m_{1,n}+1$ changes the logarithmic exponent only by $o(1)$. Hence
\[
\widetilde S_n\weakto (G,J).
\]
Moreover,
\[
n^{-1/2}\Pi_G(\mu_1-\mu_0)\to 0,
\qquad
\Pi_J(\mu_1-\mu_0)=\Delta.
\]
Therefore
\[
Q_n\weakto \mathcal L(G,J+\Delta)=Q_\infty.
\]
If $D_0\cap D_1\neq\varnothing$, then Step~1 gives $\Delta=0$, so the same weak limit appears under both
hypotheses. This proves part~(i).
\end{proof}

\section{Proof of part (ii): projected TV / Le Cam convergence}
\label{sec:proof-ii}

\begin{proof}[Proof of Theorem~\ref{thm:general-k-dominant-LK}(ii)]
Define the projected statistic
\[
S_n^J:=\Pi_J\widehat H_n,
\]
and let $P_n^J,Q_n^J$ denote its laws under $T_{n,k_n}$ and $T_{n,k_n+1}$, respectively.

\textbf{Step 1: build the quotient alphabet and the pushed-forward array.}
Let
\[
\mathcal Y_J:=\Pi_J\bigl(\{e_y:y\in\mathcal Y\}\bigr)\subset M^\perp,
\]
and for each $b\in\{0,1\}$ let $\widetilde W_b^{(n)}$ be the pushforward of $W_b^{(n)}$ under the map
$y\mapsto \Pi_J e_y$. By Lemma~\ref{lem:new-projected-array} in Appendix~\ref{app:technical-lemmas},
$\widetilde W_b^{(n)}$ has dominant quotient atom $m_b=\Pi_J\mu_b$, and for every
$w\in \mathcal Y_J\setminus\{m_b\}$,
\[
n\,\widetilde W_b^{(n)}(w)\to \widetilde \alpha_b(w):=
\sum_{\substack{y\notin D_b\\ \Pi_J e_y=w}}\alpha_b(y).
\]

\medskip
\textbf{Step 2: separate the disjoint and overlap quotient geometries.}
By Lemma~\ref{lem:new-geometry}, there are exactly two possibilities.
If $D_0\cap D_1=\varnothing$, then $m_0\neq m_1$. If $D_0\cap D_1\neq\varnothing$, then $m_0=m_1$ and
$\Delta=0$. This answers the subtlety left open in the proof skeleton: the quotient-alphabet construction is
literally Part~II, Theorem~5.8 applied to the projected array only in the disjoint case. In the overlap case one
must rerun the projected Poisson approximation with a \emph{common} dominant quotient atom.

\medskip
\textbf{Step 3: the disjoint quotient case $m_0\neq m_1$.}
Assume first that $D_0\cap D_1=\varnothing$. Then Lemma~\ref{lem:new-projected-array} shows that the pushed-
forward array $\widetilde W^{(n)}$ lies in the sparse-error critical regime of Part~II, Definition~5.1, on the
finite alphabet $\mathcal Y_J$, with dominant outputs $m_0,m_1$ and rare intensities $\widetilde\alpha_b$.
Moreover, the quantitative hypothesis \eqref{eq:quantitative-critical-regime} transfers to the projected array:
for every $w\neq m_b$,
\[
\bigl|n\widetilde W_b^{(n)}(w)-\widetilde\alpha_b(w)\bigr|
\le \sum_{\substack{y\notin D_b\\ \Pi_J e_y=w}} \bigl|nW_b^{(n)}(y)-\alpha_b(y)\bigr|
\le \frac{C}{n},
\]
because each fiber of $y\mapsto \Pi_J e_y$ is finite; and Lemma~\ref{lem:new-projected-array} gives
$\widetilde W_b^{(n)}(m_b)=1-O(n^{-1})$. Thus the projected array satisfies the quantitative condition required
for Part~II, Theorem~5.8. Concretely, that theorem applies to a finite alphabet with two distinct dominant atoms
$m_0\neq m_1$, all remaining coordinates of order $n^{-1}$ with $O(n^{-1})$ parameter control, and any composition
sequence $\pi_n\to\pi$; it yields an $O(n^{-1})$ total-variation approximation of the centered projected histogram
by the corresponding compound-Poisson / Poisson-shift limit experiment. Let $\widetilde H_n$ be the associated
centered quotient histogram under $T_{n,k_n}$, and let
\[
L:\mathbb Z^{\mathcal Y_J}\to M^\perp,
\qquad L(h):=\sum_{w\in\mathcal Y_J} h(w)\,w.
\]
By construction,
\[
S_n^J=L(\widetilde H_n).
\]
Part~II, Theorem~5.8 applied to $\widetilde W^{(n)}$ gives
\[
\TV\bigl(\mathcal L(\widetilde H_n),\widetilde P_\infty\bigr)+
\TV\bigl(\mathcal L(\widetilde H_n^{\mathrm{alt}}),\widetilde Q_\infty\bigr)
\le \frac{C}{n},
\]
with $\widetilde P_\infty,\widetilde Q_\infty$ the quotient Poisson-shift limits on $\mathbb Z^{\mathcal Y_J}$. Applying
Lemma~\ref{lem:tv-contraction-main} to the measurable map $L$ yields
\[
\TV(P_n^J,P_\infty^J)+\TV(Q_n^J,Q_\infty^J)\le \frac{C}{n},
\]
where $P_\infty^J=\mathcal L(J)$ and $Q_\infty^J=\mathcal L(J+\Delta)$ after grouping equal projected rare
atoms.

\medskip
\textbf{Step 4: the overlap quotient case $m_0=m_1$.}
Assume now that $D_0\cap D_1\neq\varnothing$, so that $m_0=m_1=:m_\star$ and $\Delta=0$.
Let
\[
\mathcal Y_J^\star:=\mathcal Y_J\setminus\{m_\star\}.
\]
For $b\in\{0,1\}$ and under $T_{n,k_n}$, define the projected rare-count vector
\[
D_{b,n}:=\bigl(\widetilde N_{b,n}(w)\bigr)_{w\in \mathcal Y_J^\star},
\]
where $\widetilde N_{0,n}\sim \Mult(m_{0,n},\widetilde W_0^{(n)})$ and
$\widetilde N_{1,n}\sim \Mult(m_{1,n},\widetilde W_1^{(n)})$ are independent quotient histograms for the two groups.
Because all dominant outputs map to $m_\star$, the projected centered statistic is the deterministic linear image
\[
S_n^J=\sum_{w\in\mathcal Y_J^\star} \bigl(D_{0,n}(w)+D_{1,n}(w)\bigr)(w-m_\star).
\]
For each $b$, the total projected rare mass
\[
\widetilde p_{b,n}:=\sum_{w\in\mathcal Y_J^\star}\widetilde W_b^{(n)}(w)
\]
is $O(n^{-1})$ by Lemma~\ref{lem:new-projected-array}. The explicit Poisson approximation for multinomials on
rare categories from Part~II, Appendix~A (the same estimate used in the proof of Part~II, Theorem~5.8), gives
independent Poisson vectors $\widetilde U_{b,n}$ on $\mathcal Y_J^\star$ with means
$m_{b,n}\widetilde W_b^{(n)}(w)$ such that
\[
\TV\bigl(\mathcal L(D_{b,n}),\mathcal L(\widetilde U_{b,n})\bigr)
\le m_{b,n}\widetilde p_{b,n}\bigl(1-e^{-\widetilde p_{b,n}}\bigr)
\le \frac{C}{n}.
\]
Since the dimensions are finite and assumption~\eqref{eq:quantitative-critical-regime} controls the means at rate
$O(n^{-1})$, perturbing the Poisson means to their limits
\[
(1-\pi)\widetilde\alpha_0(w),\qquad \pi\widetilde\alpha_1(w),
\qquad w\in\mathcal Y_J^\star,
\]
changes the law by at most $C/n$ in total variation (coordinatewise Poisson perturbation plus
Lemma~\ref{lem:tv-tensorization-main}). The sum of the two independent limiting Poisson vectors has exactly the law
of the compound-Poisson jump field $J$ on the quotient space. Applying the deterministic linear map
\[
L_\star(h):=\sum_{w\in\mathcal Y_J^\star} h(w)(w-m_\star)
\]
and using Lemma~\ref{lem:tv-contraction-main} yields
\[
\TV(P_n^J,\mathcal L(J))\le \frac{C}{n}.
\]
The same argument applies under the neighboring alternative. Since $\Delta=0$, the projected limits under both
hypotheses coincide and equal $\mathcal L(J)$. Hence
\[
\TV(P_n^J,P_\infty^J)+\TV(Q_n^J,Q_\infty^J)\le \frac{C}{n}
\]
in the overlap case as well.

\medskip
\textbf{Step 5: identify the projected limit.}
In the disjoint case the limit is the Poisson-shift pair from Step~3; in the overlap case it is the common
compound-Poisson law from Step~4. In either case the laws are exactly $P_\infty^J=\mathcal L(J)$ and
$Q_\infty^J=\mathcal L(J+\Delta)$ from the statement.

\medskip
\textbf{Step 6: conclude Le Cam convergence.}
The two projected experiments live on the same measurable space $M^\perp$, so
Lemma~\ref{lem:lecam-tv-bound-main} gives
\[
\LC\bigl((P_n^J,Q_n^J),(P_\infty^J,Q_\infty^J)\bigr)
\le
\max\{\TV(P_n^J,P_\infty^J),\TV(Q_n^J,Q_\infty^J)\}
\le \frac{C}{n}.
\]
This proves part~(ii).
\end{proof}

\section{\texorpdfstring{Proof of part (iii): \privacycurve\ convergence}{Proof of part (iii): privacy-curve convergence}}
\label{sec:proof-iii}

Write again
\[
G_n:=n^{-1/2}\Pi_G\widehat H_n,
\qquad
Z_n:=\Pi_J\widehat H_n.
\]
For the regular conditional laws of $G_n$ given $Z_n=z$ under $T_{n,k_n}$ and $T_{n,k_n+1}$ we write,
respectively,
\[
P_{n,z}^{G\mid Z}
\qquad\text{and}\qquad
Q_{n,z}^{G\mid Z}.
\]


The proof has five moving parts. In the disjoint interior regime we first recover the dominant totals from the
quotient statistic, then smooth the conditional dominant blocks by replacing a small number of foreign messages
with native dominant draws, then attach the resulting common conditional factor to the projected quotient
experiment and reduce the \privacycurve\ to the projected one by Blackwell equivalence. After that, the overlap
regime is handled directly by a one-user replacement argument. Finally, the boundary regimes are treated
separately: the weak boundary still admits conditional smoothing, while the strong boundary requires an
additional structural hypothesis and is followed by a counterexample showing why no fully general theorem can hold.

\subsection*{Interior disjoint case: Steps 1--8}

\begin{proof}[Proof of Theorem~\ref{thm:general-k-dominant-LK}(iii): interior disjoint case]
Assume throughout this subsection that $\pi\in(0,1)$ and $D_0\cap D_1=\varnothing$.

\textbf{Step 1: recover the dominant totals from the projected statistic.}
Lemma~\ref{lem:new-measurable-totals} from Appendix~\ref{app:technical-lemmas} shows that for each $b\in\{0,1\}$
there exists a measurable map $\tau_{b,n}:M^\perp\to\N$ such that
\[
N_{n,k_n}(D_b)=\tau_{b,n}(Z_n)\qquad\text{a.s. under both hypotheses.}
\]
For later use write
\[
T_{b,n}(z):=\tau_{b,n}(z).
\]
This is precisely what replaces the pair-total measurability of Part~II, Appendix~B.

\medskip
\textbf{Step 2: define the native conditional kernel.}
For $y\in D_b$ define the normalized dominant law
\[
\vartheta_{b,n}(y):=\frac{W_b^{(n)}(y)}{W_b^{(n)}(D_b)}.
\]
Given $z$ in the support of $Z_n$, let $X_{0,n}^\circ\sim \Mult(T_{0,n}(z),\vartheta_{0,n})$ and
$X_{1,n}^\circ\sim \Mult(T_{1,n}(z),\vartheta_{1,n})$ be independent. Define
\[
\Psi_{n,z}(x_0,x_1)
:=
 n^{-1/2}\Pi_G\Bigl(
 \sum_{y\in D_0}(x_0(y)-T_{0,n}(z)\vartheta_{0,n}(y))e_y
 +\sum_{y\in D_1}(x_1(y)-T_{1,n}(z)\vartheta_{1,n}(y))e_y
 \Bigr),
\]
and let $R_{n,z}$ denote the law of $\Psi_{n,z}(X_{0,n}^\circ,X_{1,n}^\circ)$. This is exactly the common
finite-$n$ conditional factor from Lemma~\ref{lem:new-conditional-smoothing}.

\medskip
\textbf{Step 3: condition on the refined rare configuration.}
Fix one of the two hypotheses. Write $(m_{0,n}^\star,m_{1,n}^\star)$ for the actual group sizes under that
hypothesis:
\[
(m_{0,n}^\star,m_{1,n}^\star)=
\begin{cases}
(m_{0,n},m_{1,n}), & \text{under }T_{n,k_n},\\
(m_{0,n}-1,m_{1,n}+1), & \text{under }T_{n,k_n+1}.
\end{cases}
\]
Refine the conditioning by revealing: (a) which users in each group produce outputs outside their native dominant
set; (b) among those rare users, which ones land in the opposite dominant set and which ones land elsewhere; and
(c) all rare outputs outside $D_0\cup D_1$. Under this refined conditioning, the dominant blocks are independent and
have the form
\[
X_{b,n}=S_{b,n}+\sum_{r=1}^{A_{b,n}}\eta_{b,r},
\]
where
\[
S_{b,n}\sim \Mult(m_{b,n}^\star-L_{b,n},\vartheta_{b,n})
\]
is the native dominant contribution, $L_{b,n}$ is the number of group-$b$ users leaving $D_b$, $A_{b,n}$ is the
number of users from the opposite group landing in $D_b$, and the $\eta_{b,r}$ are i.i.d. categorical random vectors
in $\{e_y:y\in D_b\}$ with some cross law $\rho_{b,n}$ determined by the opposite group. The surrogate kernel
$R_{n,z}$ replaces each $\eta_{b,r}$ by an independent native draw $\eta_{b,r}^\circ\sim \Cat(\vartheta_{b,n})$.

\medskip
\textbf{Step 4: telescoping TV bound inside each dominant block.}
Fix $b\in\{0,1\}$ and condition on the refined rare configuration from Step~3. Let
\[
Y_{b,n}^{(0)}:=S_{b,n}+\sum_{r=1}^{A_{b,n}}\eta_{b,r},
\qquad
Y_{b,n}^{(A_{b,n})}:=S_{b,n}+\sum_{r=1}^{A_{b,n}}\eta_{b,r}^\circ,
\]
and for $0\le s\le A_{b,n}$ define the interpolating sums
\[
Y_{b,n}^{(s)}:=S_{b,n}+\sum_{r=1}^{s}\eta_{b,r}^\circ+\sum_{r=s+1}^{A_{b,n}}\eta_{b,r}.
\]
By the triangle inequality,
\[
\TV\bigl(\mathcal L(Y_{b,n}^{(0)}),\mathcal L(Y_{b,n}^{(A_{b,n})})\bigr)
\le \sum_{s=0}^{A_{b,n}-1}
\TV\bigl(\mathcal L(Y_{b,n}^{(s)}),\mathcal L(Y_{b,n}^{(s+1)})\bigr).
\]
Now condition additionally on all variables except the pair $(\eta_{b,s+1},\eta_{b,s+1}^\circ)$. The common part is a
random vector independent of this pair. Since $p_b(y)>0$ for every $y\in D_b$ and $W_b^{(n)}(y)=p_b(y)+O(n^{-1})$, the normalized dominant law satisfies $\vartheta_{b,n}(y)\ge p_b(y)/2>0$ for all sufficiently large $n$, so the hypothesis $\vartheta_{\min}>0$ of Lemma~\ref{lem:new-edge-shift} is met with a constant depending only on $p_b$. By that lemma, the shift of a
multinomial block by one category-vector to another costs at most $C_b/\sqrt{m_{b,n}^\star-L_{b,n}+1}$ in total
variation, uniformly over the two categories. Averaging over the pair
$(\eta_{b,s+1},\eta_{b,s+1}^\circ)$ therefore gives
\[
\TV\bigl(\mathcal L(Y_{b,n}^{(s)}),\mathcal L(Y_{b,n}^{(s+1)})\bigr)
\le \frac{C_b}{\sqrt{m_{b,n}^\star-L_{b,n}+1}}.
\]
Summing over $s$ yields
\[
\TV\bigl(\mathcal L(Y_{b,n}^{(0)}),\mathcal L(Y_{b,n}^{(A_{b,n})})\bigr)
\le \frac{C_bA_{b,n}}{\sqrt{m_{b,n}^\star-L_{b,n}+1}}.
\]
Since the two dominant blocks are conditionally independent and the map $\Psi_{n,z}$ is measurable, we obtain
\[
\TV\bigl(\mathcal L(G_n\mid Z_n,\mathcal F_n),R_{n,Z_n}\bigr)
\le
\frac{C_0A_{0,n}}{\sqrt{m_{0,n}^\star-L_{0,n}+1}}
+
\frac{C_1A_{1,n}}{\sqrt{m_{1,n}^\star-L_{1,n}+1}},
\]
where $\mathcal F_n$ denotes the refined rare sigma-field.

\medskip
\textbf{Step 5: average over the rare configuration.}
Because $\pi\in(0,1)$, both actual group sizes are of order $n$ under both hypotheses. More precisely, there exists
$\kappa>0$ such that for all large $n$,
\[
m_{0,n}^\star\ge 3\kappa n,
\qquad
m_{1,n}^\star\ge 3\kappa n.
\]
The total numbers of rare outputs and cross messages have bounded means, because the rare probabilities are
$O(n^{-1})$ uniformly and the alphabet is finite. Hence
\[
\E[A_{b,n}]=O(1),\qquad \E[L_{b,n}]=O(1),\qquad b\in\{0,1\}.
\]
On the event $\{L_{b,n}\le \kappa n\}$ we have
\[
\frac{A_{b,n}}{\sqrt{m_{b,n}^\star-L_{b,n}+1}}
\le \frac{A_{b,n}}{\sqrt{2\kappa n}},
\]
while on the complementary event the same quantity is bounded by $A_{b,n}$. Let
\[
R_n:=L_{0,n}+L_{1,n}.
\]
Then $A_{b,n}\le R_n$, $L_{b,n}\le R_n$, and the rare-count estimates give $\E[R_n^2]=O(1)$. Therefore
\[
\E\Bigl[A_{b,n}\1\{L_{b,n}>\kappa n\}\Bigr]
\le
\E\Bigl[R_n\1\{R_n>\kappa n\}\Bigr]
\le
\frac{\E[R_n^2]}{\kappa n}
=O(n^{-1}).
\]
Therefore
\[
\E\Bigl[\frac{A_{b,n}}{\sqrt{m_{b,n}^\star-L_{b,n}+1}}\Bigr]=O(n^{-1/2}).
\]
Integrating the Step~4 bound over the refined rare configuration and then over $Z_n$ proves
Lemma~\ref{lem:new-conditional-smoothing}:
\[
\int \TV\bigl(P_{n,z}^{G\mid Z},R_{n,z}\bigr)\,P_n^J(dz)
+
\int \TV\bigl(Q_{n,z}^{G\mid Z},R_{n,z}\bigr)\,Q_n^J(dz)
\le \frac{C_{\mathrm{int}}}{\sqrt n}.
\]

\medskip
\textbf{Step 6: build the auxiliary experiment with common conditional factor.}
Define probability laws on $M\times M^\perp$ by
\[
\widetilde P_n(dg,dz):=P_n^J(dz)R_{n,z}(dg),
\qquad
\widetilde Q_n(dg,dz):=Q_n^J(dz)R_{n,z}(dg).
\]
Let $K((g,z),B):=\1_B(z)$ be the projection kernel onto $M^\perp$, and let
\[
L(z,A):=\int \1_A(g,z)\,R_{n,z}(dg)
\]
be the enrichment kernel. Then
\[
\widetilde P_nK=P_n^J,
\qquad \widetilde Q_nK=Q_n^J,
\qquad P_n^J L=\widetilde P_n,
\qquad Q_n^J L=\widetilde Q_n.
\]
Thus the projected experiment $(P_n^J,Q_n^J)$ and the auxiliary experiment
$(\widetilde P_n,\widetilde Q_n)$ are linked by Markov kernels in both directions. Lemma~\ref{lem:privacy-blackwell-equivalence}
therefore gives
\[
\delta_{\widetilde Q_n\|\widetilde P_n}(\varepsilon)=\delta_{Q_n^J\|P_n^J}(\varepsilon).
\]

\medskip
\textbf{Step 7: compare the full experiment to the auxiliary experiment.}
By Lemma~\ref{lem:privacy-stability-main},
\begin{align*}
\bigl|\delta_{Q_n\|P_n}(\varepsilon)-\delta_{Q_n^J\|P_n^J}(\varepsilon)\bigr|
&=
\bigl|\delta_{Q_n\|P_n}(\varepsilon)-\delta_{\widetilde Q_n\|\widetilde P_n}(\varepsilon)\bigr| \\
&\le
\TV(Q_n,\widetilde Q_n)+e^\varepsilon\TV(P_n,\widetilde P_n).
\end{align*}
Disintegrating $P_n$ and $Q_n$ with respect to $Z_n$ and using the defining property of $\widetilde P_n$ and
$\widetilde Q_n$, the two total-variation terms are exactly the two integrals controlled in Step~5. Therefore
\[
\bigl|\delta_{Q_n\|P_n}(\varepsilon)-\delta_{Q_n^J\|P_n^J}(\varepsilon)\bigr|
\le \frac{C_{\mathrm{int}}(1+e^\varepsilon)}{\sqrt n}.
\]

\medskip
\textbf{Step 8: pass to the projected limit and remove the common Gaussian factor.}
By part~(ii) and Lemma~\ref{lem:privacy-stability-main},
\[
\delta_{Q_n^J\|P_n^J}(\varepsilon)\to \delta_{Q_\infty^J\|P_\infty^J}(\varepsilon)
=\delta_{\mathcal L(J+\Delta)\|\mathcal L(J)}(\varepsilon)
\]
for every fixed $\varepsilon\ge 0$. Since $P_\infty=\mathcal L(G,J)$ and $Q_\infty=\mathcal L(G,J+\Delta)$, and $G$ is
common and independent under both hypotheses, Lemma~\ref{lem:common-independent-factor-main} gives
\[
\delta_{Q_\infty\|P_\infty}(\varepsilon)
=
\delta_{\mathcal L(J+\Delta)\|\mathcal L(J)}(\varepsilon).
\]
Combining these two displays with Step~7 proves part~(iii) in the interior disjoint case.
\end{proof}

\subsection*{Interior overlap case: Steps 9--11}

\begin{proof}[Proof of Theorem~\ref{thm:general-k-dominant-LK}(iii): interior overlap case]
Assume $\pi\in(0,1)$ and $D_0\cap D_1\neq\varnothing$.

\textbf{Step 9: identify the limiting privacy problem.}
By Theorem~\ref{thm:general-k-dominant-LK}(i),
\[
\Delta=0,
\qquad
P_\infty=Q_\infty=\mathcal L(G,J),
\qquad
\delta_{Q_\infty\|P_\infty}(\varepsilon)\equiv 0.
\]
Thus it suffices to prove that $\delta_{Q_n\|P_n}(\varepsilon)\to 0$.

\medskip
\textbf{Step 10: prove a direct $O(n^{-1/2})$ TV comparison between the full laws.}
Choose and fix a shared dominant symbol $s\in D_0\cap D_1$. Under $T_{n,k_n}$ write the dataset as
$n-k_n-1$ common zeros, $k_n$ common ones, and one distinguished extra zero. Under $T_{n,k_n+1}$ keep the same
$n-1$ common users but replace the distinguished extra zero by a distinguished extra one. Couple all common users
identically under the two hypotheses. Let $H_n^{\mathrm{com}}$ be the common contribution of those $n-1$ users after centering
by $(n-k_n-1)\mu_0+k_n\mu_1$. Then under the two hypotheses,
\[
S_n^{(P)}=H_n^{\mathrm{com}}+\bigl(n^{-1/2}\Pi_G(e_{A_n}-\mu_0),\ \Pi_J(e_{A_n}-\mu_0)\bigr),
\]
\[
S_n^{(Q)}=H_n^{\mathrm{com}}+\bigl(n^{-1/2}\Pi_G(e_{B_n}-\mu_0),\ \Pi_J(e_{B_n}-\mu_0)\bigr),
\]
where $A_n\sim W_0^{(n)}$ and $B_n\sim W_1^{(n)}$ are independent of the common part. Rare outputs of the switched
user have probability $O(n^{-1})$ under either hypothesis, so their total contribution to the total-variation
distance is $O(n^{-1})$.

On the dominant event $\{A_n\in D_0,\ B_n\in D_1\}$, we use the triangle inequality with the common anchor $s$:
\begin{align*}
\TV\bigl(\mathcal L(S_n^{(P)}),\mathcal L(S_n^{(Q)})\bigr)
&\le
\TV\bigl(\mathcal L(H_n^{\mathrm{com}}+\Xi(A_n)),\mathcal L(H_n^{\mathrm{com}}+\Xi(s))\bigr) \\
&\qquad+
\TV\bigl(\mathcal L(H_n^{\mathrm{com}}+\Xi(s)),\mathcal L(H_n^{\mathrm{com}}+\Xi(B_n))\bigr)
+O(n^{-1}),
\end{align*}
where $\Xi(y):=(n^{-1/2}\Pi_G(e_y-\mu_0),\Pi_J(e_y-\mu_0))$. To bound the first term, condition on the entire common
$1$-group and all rare outputs from the common $0$-group. The remaining dominant $0$-group block is a multinomial
with native law $\vartheta_{0,n}$ and sample size of order $n$; adding $A_n$ versus $s$ changes it by a single edge
shift, so Lemma~\ref{lem:new-edge-shift} gives an $O(n^{-1/2})$ bound. The second term is controlled in the same way
by conditioning instead on the common $0$-group and smoothing inside the common $1$-group block. Averaging over the
conditioning variables yields
\[
\TV(P_n,Q_n)\le \frac{C}{\sqrt n},
\]
which is Lemma~\ref{lem:new-overlap-tv} from Appendix~\ref{app:technical-lemmas}.

\medskip
\textbf{Step 11: conclude \privacycurve\ convergence.}
By the definition of the privacy curve,
\[
0\le \delta_{Q_n\|P_n}(\varepsilon)\le \sup_A |Q_n(A)-P_n(A)|\le \TV(P_n,Q_n).
\]
Step~10 therefore implies
\[
\delta_{Q_n\|P_n}(\varepsilon)\le \frac{C}{\sqrt n}\to 0
=\delta_{Q_\infty\|P_\infty}(\varepsilon).
\]
This proves part~(iii) in the interior overlap case.
\end{proof}


\subsection*{Boundary case: Step 12}

\begin{proof}[Proof of Theorem~\ref{thm:general-k-dominant-LK}(iii): boundary regimes covered by the theorem]
Assume first that $\pi=0$; the case $\pi=1$ is symmetric after swapping the roles of the two groups.
Write
\[
m_{0,n}:=n-k_n,
\qquad
m_{1,n}:=k_n.
\]

\textbf{(a) Strong boundary with the regularity condition from part~\textup{(iii)(c)}.}
Assume $k_n=O(1)$.

\smallskip
\emph{Case 1: the bounded minority dominant block is absorbed by the majority block, i.e. $D_1\subseteq D_0$.}
In this case Lemma~\ref{lem:new-geometry} gives $\Delta=0$, so it suffices to show
\[
\TV(P_n,Q_n)\to 0.
\]
Fix $K<\infty$ with $k_n\le K$ for all large $n$. Couple the two neighboring datasets by using the same
$n-k_n-1$ common $0$-users and the same $k_n$ common $1$-users under both hypotheses, and by switching only one
distinguished user from input $0$ under $T_{n,k_n}$ to input $1$ under $T_{n,k_n+1}$. Reveal all outputs outside
$D_0$ of the common users and, under each hypothesis, also reveal whether the switched user leaves $D_0$.
Because $D_1\subseteq D_0$ and all rare probabilities are $O(n^{-1})$, the probability that any of the $k_n$
common $1$-users or the switched user leaves $D_0$ is $O(n^{-1})$.

Conditional on the revealed rare data and on the event that all common $1$-users and the switched user stay in
$D_0$, the remaining common $0$-user block on $D_0$ is distributed as
\[
S_n\sim \Mult(m_n',\vartheta_{0,n}),
\qquad
m_n'=n-O_{\Pbb}(1),
\]
where $\vartheta_{0,n}(y):=W_0^{(n)}(y)/W_0^{(n)}(D_0)$ and $m_n'$ equals the number of common $0$-users whose
outputs stay in $D_0$. The contribution of the common $1$-users and of the switched user consists of at most
$K+1$ additional category vectors in $D_0$ under each hypothesis. Hence, conditional on the revealed rare data and
on the values of those $O(1)$ minority outputs, both full statistics are measurable images of
\[
S_n+\sum_{r=1}^{R_n} e_{a_r}
\qquad\text{and}\qquad
S_n+\sum_{r=1}^{R_n'} e_{b_r},
\]
for some $R_n,R_n'\le K+1$ and some $a_r,b_r\in D_0$. By the multi-step part of
Lemma~\ref{lem:new-edge-shift},
\[
\TV\!\left(
\mathcal L\!\left(S_n+\sum_{r=1}^{R_n} e_{a_r}\,\middle|\,\mathcal F_n\right),
\mathcal L\!\left(S_n+\sum_{r=1}^{R_n'} e_{b_r}\,\middle|\,\mathcal F_n\right)
\right)
\le
\frac{C(R_n+R_n')}{\sqrt{m_n'+1}}
\le
\frac{C_K}{\sqrt{m_n'+1}},
\]
where $\mathcal F_n$ denotes the revealed sigma-field. Since the number of common rare outputs has bounded second
moment, the same argument as in Lemma~\ref{lem:new-conditional-smoothing} yields
\[
\E\bigl[(m_n'+1)^{-1/2}\bigr]=O(n^{-1/2}).
\]
Therefore the conditional total-variation distance is $O(n^{-1/2})$ on the event that all minority users stay in
$D_0$, while the complementary event has probability $O(n^{-1})$. Averaging proves
\[
\TV(P_n,Q_n)\le \frac{C}{\sqrt n}.
\]
Hence, for every fixed $\varepsilon\ge 0$,
\[
0\le \delta_{Q_n\|P_n}(\varepsilon)\le \TV(P_n,Q_n)\to 0
=\delta_{Q_\infty\|P_\infty}(\varepsilon).
\]

\smallskip
\emph{Case 2: the bounded minority dominant block is one-point and disjoint from the majority block, i.e.
$|D_1|=1$ and $D_1\cap D_0=\varnothing$.}
Write $D_1=\{y_\star\}$. Then $M_1=\{0\}$, and the only dominant Gaussian block is the majority $D_0$-block.
Since the case is disjoint, Lemma~\ref{lem:new-measurable-totals} applies and the dominant total
\[
T_{0,n}(z):=N_{n,k_n}(D_0)
\]
is $\sigma(Z_n)$-measurable. Define
\[
\vartheta_{0,n}(y):=\frac{W_0^{(n)}(y)}{W_0^{(n)}(D_0)},
\qquad y\in D_0,
\]
and for $z$ in the support of $Z_n$ let $X_{0,n}^\circ\sim \Mult(T_{0,n}(z),\vartheta_{0,n})$. Let $R_{n,z}$ be
the law of
\[
\Phi_{n,z}(X_{0,n}^\circ)
:=
n^{-1/2}\Pi_G\!\left(
\sum_{y\in D_0}(X_{0,n}^\circ(y)-T_{0,n}(z)\vartheta_{0,n}(y))e_y
\right).
\]
Define the auxiliary experiment
\[
\widetilde P_n(dg,dz):=P_n^J(dz)\,R_{n,z}(dg),
\qquad
\widetilde Q_n(dg,dz):=Q_n^J(dz)\,R_{n,z}(dg).
\]
By Lemma~\ref{lem:privacy-blackwell-equivalence},
\[
\delta_{\widetilde Q_n\|\widetilde P_n}(\varepsilon)=\delta_{Q_n^J\|P_n^J}(\varepsilon).
\]

It remains to compare $(P_n,Q_n)$ to $(\widetilde P_n,\widetilde Q_n)$. Work under either hypothesis and let
$m_{0,n}^\star=n-k_n+O(1)$ denote the actual number of $0$-users. Reveal all outputs outside $D_0$ and let $L_{0,n}$
be the number of $0$-users leaving $D_0$. Let $B_n$ be the number of minority users landing in $D_0$; because
$|D_1|=1$ and $D_1\cap D_0=\varnothing$, every such landing is a rare cross event and therefore
\[
\E[B_n]=O(n^{-1}).
\]
Conditional on the revealed rare data, the $D_0$-block has the form
\[
Y_n=S_n+\sum_{r=1}^{B_n}\eta_r,
\qquad
S_n\sim \Mult(m_{0,n}^\star-L_{0,n},\vartheta_{0,n}),
\]
where the $\eta_r$ are $D_0$-valued category vectors coming from the minority users that cross into $D_0$.
Replace them one by one by independent native draws $\eta_r^\circ\sim \Cat(\vartheta_{0,n})$, independent of
everything else, and write
\[
Y_n^\circ:=S_n+\sum_{r=1}^{B_n}\eta_r^\circ.
\]
Then, conditional on $Z_n=z$, the law of $Y_n^\circ$ is exactly $\Mult(T_{0,n}(z),\vartheta_{0,n})$, hence its
image under $\Phi_{n,z}$ is $R_{n,z}$. By Lemma~\ref{lem:new-edge-shift},
\[
\TV\bigl(\mathcal L(Y_n\mid \mathcal F_n),\mathcal L(Y_n^\circ\mid \mathcal F_n)\bigr)
\le \frac{CB_n}{\sqrt{m_{0,n}^\star-L_{0,n}+1}},
\]
and therefore
\[
\TV\bigl(\mathcal L(G_n\mid Z_n,\mathcal F_n),R_{n,Z_n}\bigr)
\le \frac{CB_n}{\sqrt{m_{0,n}^\star-L_{0,n}+1}}.
\]
Since $L_{0,n}$ has bounded second moment and $\E[B_n]=O(n^{-1})$, averaging yields
\[
\TV(P_n,\widetilde P_n)+\TV(Q_n,\widetilde Q_n)=o(1).
\]
By Lemma~\ref{lem:privacy-stability-main},
\[
\bigl|\delta_{Q_n\|P_n}(\varepsilon)-\delta_{Q_n^J\|P_n^J}(\varepsilon)\bigr|=o(1).
\]
Part~(ii) then gives
\[
\delta_{Q_n\|P_n}(\varepsilon)\to
\delta_{Q_\infty^J\|P_\infty^J}(\varepsilon)
=
\delta_{Q_\infty\|P_\infty}(\varepsilon).
\]

\medskip
\textbf{(b) Weak boundary: $k_n\to\infty$ and $k_n/n\to 0$.}
We first treat the disjoint case $D_0\cap D_1=\varnothing$. Define $\vartheta_{b,n}$ and the surrogate kernels
$R_{n,z}$ exactly as in Steps~2 and~6 of Section~\ref{sec:proof-iii}. Returning to the refined conditioning from
Step~3, the same telescoping argument as in Step~4 gives, under either hypothesis,
\[
\TV\bigl(\mathcal L(G_n\mid Z_n,\mathcal F_n),R_{n,Z_n}\bigr)
\le
\frac{C_0A_{0,n}}{\sqrt{m_{0,n}^\star-L_{0,n}+1}}
+
\frac{C_1A_{1,n}}{\sqrt{m_{1,n}^\star-L_{1,n}+1}},
\]
where $m_{0,n}^\star=n-k_n+O(1)$ and $m_{1,n}^\star=k_n+O(1)$. Here $A_{0,n}$ counts minority users landing in
$D_0$, while $A_{1,n}$ counts majority users landing in $D_1$. Since each user hits the opposite dominant set with
probability $O(n^{-1})$,
\[
\E[A_{0,n}]=O(k_n/n),
\qquad
\E[A_{1,n}]=O(1),
\qquad
\E[L_{0,n}]=O(1),
\qquad
\E[L_{1,n}]=O(k_n/n).
\]
Because $m_{0,n}^\star\asymp n$ and $m_{1,n}^\star\asymp k_n$, Markov's inequality gives
\[
\Pbb\bigl(L_{0,n}>\tfrac12 m_{0,n}^\star\bigr)=O(n^{-1}),
\qquad
\Pbb\bigl(L_{1,n}>\tfrac12 m_{1,n}^\star\bigr)=O(n^{-1}).
\]
Using the bounded second moments of the rare counts exactly as in Step~5, we obtain
\[
\E\Bigl[\frac{A_{0,n}}{\sqrt{m_{0,n}^\star-L_{0,n}+1}}\Bigr]=O(n^{-1/2}),
\qquad
\E\Bigl[\frac{A_{1,n}}{\sqrt{m_{1,n}^\star-L_{1,n}+1}}\Bigr]=O(k_n^{-1/2}).
\]
Therefore
\[
\int \TV\bigl(P_{n,z}^{G\mid Z},R_{n,z}\bigr)\,P_n^J(dz)
+
\int \TV\bigl(Q_{n,z}^{G\mid Z},R_{n,z}\bigr)\,Q_n^J(dz)
=O(n^{-1/2})+O(k_n^{-1/2})=o(1).
\]
Defining $(\widetilde P_n,\widetilde Q_n)$ exactly as in Step~6 and applying Step~7 gives
\[
\bigl|\delta_{Q_n\|P_n}(\varepsilon)-\delta_{Q_n^J\|P_n^J}(\varepsilon)\bigr|=o(1).
\]
At this point we do \emph{not} invoke part~\textup{(ii)}, because the quantitative hypothesis
\eqref{eq:quantitative-critical-regime} need not hold when $k_n\to\infty$ and $k_n/n\to 0$. Instead, return to the
quotient alphabet construction from Section~\ref{sec:proof-ii}. Let $\widetilde H_n$ and $\widetilde H_n^{\mathrm{alt}}$
denote the centered quotient histograms on the finite quotient alphabet $\mathcal Y_J$, and let
\[
L:\mathbb Z^{\mathcal Y_J}\to M^\perp,
\qquad
L(h):=\sum_{w\in\mathcal Y_J} h(w)\,w.
\]
In the present disjoint case the pushed-forward array still satisfies the qualitative sparse-error regime of
Part~II, Definition~5.1 on $\mathcal Y_J$, with composition parameter $\pi_n\to 0$. Hence the qualitative part of
Part~II, Theorem~5.8 yields
\[
\TV\bigl(\mathcal L(\widetilde H_n),\widetilde P_\infty\bigr)
+
\TV\bigl(\mathcal L(\widetilde H_n^{\mathrm{alt}}),\widetilde Q_\infty\bigr)
\longrightarrow 0,
\]
where $\widetilde P_\infty,\widetilde Q_\infty$ are the quotient Poisson-shift limits. Equivalently, since
$\widetilde H_n$ takes values in the discrete countable space $\mathbb Z^{\mathcal Y_J}$, one may view this as weak
convergence upgraded to total variation by pointwise convergence of the probability mass functions and Scheff\'e's
lemma. Applying the measurable map $L$ and using Lemma~\ref{lem:tv-contraction-main} gives
\[
\TV(P_n^J,P_\infty^J)+\TV(Q_n^J,Q_\infty^J)\longrightarrow 0.
\]
Therefore
\[
\delta_{Q_n\|P_n}(\varepsilon)
\to
\delta_{Q_\infty^J\|P_\infty^J}(\varepsilon)
=
\delta_{Q_\infty\|P_\infty}(\varepsilon).
\]

If instead $D_0\cap D_1\neq\varnothing$, then Lemma~\ref{lem:new-geometry} gives $\Delta=0$, so
$\delta_{Q_\infty\|P_\infty}(\varepsilon)\equiv 0$. Repeat the coupling argument from Step~10. The common
$0$-group dominant block has size $n-k_n-1\asymp n$, while the common $1$-group dominant block has size
$k_n\asymp k_n$. Smoothing inside these two blocks yields
\[
\TV(P_n,Q_n)\le \frac{C_0}{\sqrt n}+\frac{C_1}{\sqrt{k_n}}+O(n^{-1})=o(1).
\]
Hence
\[
0\le \delta_{Q_n\|P_n}(\varepsilon)\le \TV(P_n,Q_n)\to 0=\delta_{Q_\infty\|P_\infty}(\varepsilon).
\]

The case $\pi=1$ is identical after exchanging the labels of the two groups. This proves
Theorem~\ref{thm:general-k-dominant-LK}(iii) in all boundary regimes covered by the statement.
\end{proof}

\begin{proposition}[Strong-boundary obstruction]
\label{prop:strong-boundary-obstruction}
Without the additional strong-boundary regularity in Theorem~\ref{thm:general-k-dominant-LK}\textup{(iii)(c)},
the projected jump limit need not determine the full \privacycurve.
In particular, there exists an array in the regime of Definition~\ref{def:general-k-dominant} with $\pi=0$ and
$k_n\equiv 0$ such that
\[
P_\infty=Q_\infty,
\qquad
\delta_{Q_\infty\|P_\infty}(\varepsilon)\equiv 0,
\]
but
\[
\delta_{Q_n\|P_n}(\varepsilon)\ge \frac12
\qquad\text{for every }n\text{ and every }\varepsilon\ge 0.
\]
\end{proposition}

\begin{proof}
Take $\mathcal Y=\{a,b,c\}$ and define
\[
D_0=\{a,b\},
\qquad
D_1=\{b,c\},
\]
with
\[
W_0^{(n)}(a)=W_0^{(n)}(b)=\frac12,
\qquad
W_0^{(n)}(c)=0,
\]
\[
W_1^{(n)}(b)=W_1^{(n)}(c)=\frac12,
\qquad
W_1^{(n)}(a)=0.
\]
Then the hypotheses of Definition~\ref{def:general-k-dominant} hold with no rare mass, $\pi=0$, and $k_n\equiv 0$.
Because $D_0\cap D_1=\{b\}\neq\varnothing$, Lemma~\ref{lem:new-geometry} gives $\Delta=0$; moreover $\nu=0$ since
there are no rare outputs. Hence part~\textup{(i)} gives
\[
P_\infty=Q_\infty=\mathcal L(G,0),
\]
so the limiting \privacycurve\ is identically zero.

Here the overlap makes $M$ equal to the zero-sum subspace of $\R^{\{a,b,c\}}$. Under $T_{n,0}$, all outputs
lie in $\{a,b\}$, so the centered histogram already belongs to $M$ and has $c$-coordinate $0$; therefore
\[
P_n\bigl(\{(g,z):g(c)>0\}\bigr)=0.
\]
Under $T_{n,1}$ there is exactly one $1$-user. If that user outputs $c$ (probability $1/2$), then the centered
histogram still has total sum $0$ and $c$-coordinate $1$, so
\[
G_n(c)=\frac{1}{\sqrt n}>0.
\]
If the user outputs $b$ (also probability $1/2$), then the centered histogram has $c$-coordinate $0$, so
\[
G_n(c)=0.
\]
Hence
\[
Q_n\bigl(\{(g,z):g(c)>0\}\bigr)=\frac12.
\]
By the definition of the \privacycurve, for every $\varepsilon\ge 0$,
\[
\delta_{Q_n\|P_n}(\varepsilon)\ge Q_n(\{N_c>0\})-e^{\varepsilon}P_n(\{N_c>0\})=\frac12-0=\frac12.
\]
This proves the proposition.
\end{proof}

\section{The sharpness of \texorpdfstring{$O(n^{-1/2})$}{order n to minus one-half} in the hybrid experiment}
\label{sec:sharpness}

The first result of this section isolates the exact mechanism behind the $n^{-1/2}$ rate in the full hybrid
experiment. The second records the compatibility condition under which the sharper $n^{-1}$ rate returns.

\begin{theorem}[A sharp binomial-plus-Bernoulli comparison]
\label{thm:bin-ber-exact-tv}
Let $S\sim \Bin(m,p)$ with $p\in(0,1)$ and $q\in[0,1]$, and define
\[
X:=S+\mathrm{Ber}(q),
\qquad
Y:=S+\mathrm{Ber}(p),
\]
where the Bernoulli variables are independent of $S$. Then
\[
\TV\bigl(\mathcal L(X),\mathcal L(Y)\bigr)=|q-p|\,\max_k \Pbb(S=k).
\]
Consequently,
\[
\TV\bigl(\mathcal L(X),\mathcal L(Y)\bigr)
=
\frac{|q-p|}{\sqrt{2\pi mp(1-p)}}+O(m^{-1})
\qquad (m\to\infty).
\]
\end{theorem}

\begin{proof}
Let $s_k:=\Pbb(S=k)$, with the convention $s_k=0$ outside $\{0,\dots,m\}$. Then
\[
\Pbb(X=k)=(1-q)s_k+qs_{k-1},
\qquad
\Pbb(Y=k)=(1-p)s_k+ps_{k-1},
\]
so
\[
\Pbb(X=k)-\Pbb(Y=k)=(q-p)(s_{k-1}-s_k).
\]
Therefore
\[
2\TV\bigl(\mathcal L(X),\mathcal L(Y)\bigr)
=|q-p|\sum_k |s_{k-1}-s_k|.
\]
The sequence $(s_k)_k$ is binomial and hence unimodal. For any nonnegative unimodal sequence with finite support
and boundary values $0$ at both ends,
\[
\sum_k |s_{k-1}-s_k|=2\max_k s_k.
\]
Indeed, if $m_0$ is a mode, then the sequence is increasing up to $m_0$ and decreasing afterwards, so the total
variation of the sequence equals the rise from $0$ to $s_{m_0}$ plus the fall back to $0$. This proves the exact
formula.

For the asymptotic expansion, let $k_m$ be a mode of the binomial distribution. Stirling's formula yields the local
central limit estimate
\[
\Pbb(S=k_m)=\frac{1}{\sqrt{2\pi mp(1-p)}}+O(m^{-1}),
\]
see, for example, Petrov~\cite[Chapter~VII]{Pet75}. Substituting this into the exact formula proves the second
display.
\end{proof}

\begin{theorem}[Sharpness of the $n^{-1/2}$ hybrid rate]
\label{thm:sharpness-hybrid}
There exist two-dominant disjoint shuffle arrays satisfying the assumptions of Part~II, Proposition~5.4 and the
interior-composition hypothesis $\pi_n\to\pi\in(0,1)$ for which the auxiliary $O(n^{-1/2})$ comparison bound used in
Appendix~B of Part~II and in the proof of Theorem~\ref{thm:general-k-dominant-LK}(iii) is sharp. More precisely, one can construct such an array and a
corresponding auxiliary common-factor experiment $(\widetilde P_n,\widetilde Q_n)$ such that
\[
\TV(P_n,\widetilde P_n)+\TV(Q_n,\widetilde Q_n)\asymp n^{-1/2}.
\]
Consequently, the auxiliary $O(n^{-1/2})$ comparison step underlying Theorem~\ref{thm:general-k-dominant-LK}(iii)
cannot in general be improved to $o(n^{-1/2})$.
\end{theorem}

\begin{proof}
Take $\mathcal Y=\{a,b,c,d\}$, $D_0=\{a,b\}$, $D_1=\{c,d\}$, fix $p,q\in(0,1)$ with $q\neq p$, and fix
$\lambda>0$. Define the local randomizers by
\[
W_0^{(n)}(a)=p,\qquad W_0^{(n)}(b)=1-p,
\qquad W_0^{(n)}(c)=W_0^{(n)}(d)=0,
\]
\[
W_1^{(n)}(c)=p\Bigl(1-\frac{\lambda}{n}\Bigr),
\qquad
W_1^{(n)}(d)=(1-p)\Bigl(1-\frac{\lambda}{n}\Bigr),
\]
\[
W_1^{(n)}(a)=\frac{\lambda q}{n},
\qquad
W_1^{(n)}(b)=\frac{\lambda(1-q)}{n}.
\]
This is a legitimate two-dominant disjoint array. The only rare event is that a $1$-user lands in $D_0$, and the
conditional split law inside $D_0$ is $(q,1-q)$ rather than the native $0$-split $(p,1-p)$.

Let $A_n$ be the number of $1$-users landing in $D_0$. Because there are no other rare transitions in this
construction, the event
\[
E_n:=\{A_n=1\}
\]
is exactly the event that the projected jump block records a single cross message into $D_0$. Under either
hypothesis and under the auxiliary common-factor experiment, $A_n\weakto \Poi(\pi\lambda)$, so $E_n$ has
probability bounded away from $0$ and $1$; moreover $E_n$ is measurable with respect to the projected statistic and
therefore has the same probability under $P_n$ and $\widetilde P_n$. On $E_n$, the $D_0$-block of the true
experiment has the form
\[
S_n+\mathrm{Ber}(q),
\qquad S_n\sim \Bin(m_n,p),
\]
with $m_n\sim (1-\pi)n$, while the corresponding $D_0$-block in the auxiliary experiment has the form
\[
S_n+\mathrm{Ber}(p).
\]
All remaining coordinates agree under the two conditional laws on $E_n$. Therefore Theorem~\ref{thm:bin-ber-exact-tv}
yields
\[
\TV\bigl(P_n(\cdot\mid E_n),\widetilde P_n(\cdot\mid E_n)\bigr)
=|q-p|\max_k \Pbb(\Bin(m_n,p)=k)
=\frac{c_0}{\sqrt n}+O(n^{-1})
\]
for some $c_0>0$. Since $P_n(E_n)=\widetilde P_n(E_n)\to \eta\in(0,1)$, the elementary lower bound
\[
\TV(\mu,\nu)\ge \mu(E)\,\TV(\mu(\cdot\mid E),\nu(\cdot\mid E))-
|\mu(E)-\nu(E)|
\]
gives
\[
\TV(P_n,\widetilde P_n)\ge \frac{c}{\sqrt n}
\]
for all large $n$, with $c>0$. The upper bound $\TV(P_n,\widetilde P_n)=O(n^{-1/2})$ is exactly the content of
Part~II, Appendix~B; equivalently, it is the comparison estimate proved in Steps~5--7 of the proof of
Theorem~\ref{thm:general-k-dominant-LK}(iii). The same argument applies to the neighboring
law $Q_n$, proving the displayed total-variation sharpness. Since part~(ii) has projected Le Cam error $O(n^{-1})$,
the $n^{-1/2}$ bottleneck in the full hybrid comparison comes from the passage from $(P_n,Q_n)$ to
$(\widetilde P_n,\widetilde Q_n)$, and the preceding lower bound shows that this step is sharp in general.
\end{proof}

\begin{theorem}[Compatibility restores the $n^{-1}$ rate]
\label{thm:compatibility-rate}
Assume the disjoint interior setting of Theorem~\ref{thm:general-k-dominant-LK}(iii). For $b\in\{0,1\}$ and
all $n$ such that $W_{1-b}^{(n)}(D_b)>0$, let $\rho_{b,n}$ denote the conditional cross split law on $D_b$
induced by outputs from the opposite group:
\[
\rho_{0,n}(y):=\frac{W_1^{(n)}(y)}{W_1^{(n)}(D_0)},\quad y\in D_0,
\qquad
\rho_{1,n}(y):=\frac{W_0^{(n)}(y)}{W_0^{(n)}(D_1)},\quad y\in D_1.
\]
When $W_{1-b}^{(n)}(D_b)=0$, the corresponding cross block is absent and no compatibility condition is imposed.
Suppose, in addition to \eqref{eq:quantitative-critical-regime}, that
\begin{equation}
\TV(\rho_{0,n},\vartheta_{0,n})+\TV(\rho_{1,n},\vartheta_{1,n})\le \frac{C_{\mathrm{comp}}}{n}
\label{eq:compatibility-condition}
\end{equation}
for all sufficiently large $n$ for which both denominators are positive. Then the auxiliary common-factor experiment
of Section~\ref{sec:proof-iii}
satisfies
\[
\TV(P_n,\widetilde P_n)+\TV(Q_n,\widetilde Q_n)\le \frac{C}{n}
\]
for all sufficiently large $n$, and therefore
\[
\bigl|\delta_{Q_n\|P_n}(\varepsilon)-\delta_{Q_n^J\|P_n^J}(\varepsilon)\bigr|
\le \frac{C(1+e^\varepsilon)}{n}.
\]
\end{theorem}

\begin{proof}
Return to the conditional representation from Section~\ref{sec:proof-iii}. Given the refined rare configuration,
there is nothing to couple in block $b$ when $W_{1-b}^{(n)}(D_b)=0$, because then $A_{b,n}=0$ almost surely.
Otherwise the cross messages inside $D_b$ are i.i.d. with law $\rho_{b,n}$, while the surrogate kernel replaces
them by i.i.d. native draws with law $\vartheta_{b,n}$. Couple each such pair optimally. Then the probability
that a single pair mismatches is exactly $\TV(\rho_{b,n},\vartheta_{b,n})\le C_{\mathrm{comp}}/n$. Since the number
of cross messages $A_{b,n}$ has bounded expectation uniformly in $n$, the expected number of mismatches in block
$b$ is $O(n^{-1})$. Whenever all pairs match, the true and surrogate dominant blocks coincide exactly. Therefore,
conditioning on the refined rare configuration and using the coupling bound for total variation,
\[
\TV\bigl(\mathcal L(G_n\mid Z_n,\mathcal F_n),R_{n,Z_n}\bigr)
\le C\sum_{b=0}^1 A_{b,n}\,\TV(\rho_{b,n},\vartheta_{b,n}).
\]
Taking expectations and using $\E A_{b,n}=O(1)$ yields an $O(n^{-1})$ bound on the integrated conditional TV error.
Repeating the argument from Steps~6--7 in Section~\ref{sec:proof-iii} gives
\[
\TV(P_n,\widetilde P_n)+\TV(Q_n,\widetilde Q_n)\le \frac{C}{n}.
\]
The \privacycurve\ estimate follows from Lemma~\ref{lem:privacy-stability-main}.
\end{proof}

\begin{remark}[Why Theorems~\ref{thm:sharpness-hybrid} and \ref{thm:compatibility-rate} are consistent]
In the sharpness construction of Theorem~\ref{thm:sharpness-hybrid}, the only nontrivial cross block is inside
$D_0$, where the cross split law is $\rho_{0,n}=(q,1-q)$ while the native dominant law is
$\vartheta_{0,n}=(p,1-p)$. Hence
\[
\TV(\rho_{0,n},\vartheta_{0,n})=|q-p|,
\]
independently of $n$, whereas the second cross block is absent. Thus the compatibility condition
\eqref{eq:compatibility-condition} fails at order one in that example. This is exactly why
Theorem~\ref{thm:sharpness-hybrid} attains the generic $n^{-1/2}$ rate, while
Theorem~\ref{thm:compatibility-rate} restores the sharper $n^{-1}$ rate only under an additional asymptotic
matching assumption on the cross and native split laws.
\end{remark}

\section{\texorpdfstring{Boundary Berry--Esseen: Poisson/Skellam $\to$ Gaussian as $c\downarrow 0$}{Boundary Berry--Esseen: Poisson/Skellam to Gaussian as c downarrow 0}}
\label{sec:boundary-BE}

Let
\[
\mathcal E_c^{\mathrm{Poi}}:=\bigl(\Poi(c^{-2}),\ 1+\Poi(c^{-2})\bigr)
\]
denote the critical Poisson-shift experiment from Part~II. The next theorem compares it to the Gaussian shift
experiment
\[
\mathcal E_c^{\mathrm G}:=\bigl(N(0,1),\ N(c,1)\bigr)
\]
at the boundary $c\downarrow 0$.

\begin{theorem}[Boundary Berry--Esseen for the Poisson-shift experiment]
\label{thm:boundary-BE-poisson}
There exist constants $c_0>0$ and $C<\infty$ such that for every $c\in(0,c_0]$,
\[
\LC\bigl(\mathcal E_c^{\mathrm{Poi}},\mathcal E_c^{\mathrm G}\bigr)\le Cc.
\]
More concretely, after the rescaling
\[
Z_c:=c\bigl(J_c-c^{-2}\bigr),\qquad J_c\sim\Poi(c^{-2}),
\]
the experiment $\mathcal E_c^{\mathrm{Poi}}$ is equivalent to the lattice experiment
\[
\widehat{\mathcal E}_c^{\mathrm{Poi}}:=\bigl(\mathcal L(Z_c),\ \mathcal L(Z_c+c)\bigr),
\]
and there exist Markov kernels $K_c$ and $L_c$ such that
\[
\TV\bigl(K_c\mathcal L(Z_c),N(0,1)\bigr)+\TV\bigl(K_c\mathcal L(Z_c+c),N(c,1)\bigr)\le Cc,
\]
\[
\TV\bigl(L_cN(0,1),\mathcal L(Z_c)\bigr)+\TV\bigl(L_cN(c,1),\mathcal L(Z_c+c)\bigr)\le Cc.
\]
\end{theorem}

\begin{proof}
Set $\lambda:=c^{-2}$ and let
\[
z_{k,c}:=c(k-\lambda),\qquad k\in\mathbb Z.
\]
The mapping $k\mapsto z_{k,c}$ is bijective, so $\mathcal E_c^{\mathrm{Poi}}$ is exactly equivalent to the lattice
experiment $\widehat{\mathcal E}_c^{\mathrm{Poi}}$.

Let $\varphi$ be the standard normal density and write $I_{k,c}:=[z_{k,c}-c/2,z_{k,c}+c/2)$.
We claim that
\begin{equation}
\Pbb(Z_c=z_{k,c})
=
c\varphi(z_{k,c})+O\!\left(c^2(1+|z_{k,c}|)^{-3}\right)
\label{eq:poisson-local-weighted}
\end{equation}
uniformly in $k$ as $c\downarrow 0$.

Indeed, let $p_{k,\lambda}:=\Pbb(\Poi(\lambda)=k)=e^{-\lambda}\lambda^k/k!$ with $\lambda=c^{-2}$ and
$u:=z_{k,c}=c(k-\lambda)$. For $k<0$ the left-hand side is zero, so the claimed bound is immediate; hence assume
$k\ge 0$.

If $|u|\le c^{-1/3}$, then $k=\lambda+u/c=\lambda(1+uc)$ with $|uc|\le c^{2/3}$. Stirling's formula gives
\[
p_{k,\lambda}
=
\frac{1}{\sqrt{2\pi k}}
\exp\!\bigl(-\lambda+k-k\log(k/\lambda)\bigr)\bigl(1+O(k^{-1})\bigr).
\]
Since
\[
(1+t)\log(1+t)-t=\frac{t^2}{2}+O(|t|^3)
\qquad (|t|\le 1/2),
\]
we obtain
\[
-\lambda+k-k\log(k/\lambda)
=
-\lambda\bigl((1+uc)\log(1+uc)-uc\bigr)
=
-\frac{u^2}{2}+O\!\bigl(c|u|^3\bigr),
\]
while
\[
k^{-1/2}=c\bigl(1+O(c(1+|u|))\bigr).
\]
Hence, uniformly on $|u|\le c^{-1/3}$,
\[
p_{k,\lambda}
=
c\,\varphi(u)
\exp\!\bigl(O(c|u|^3)\bigr)
\bigl(1+O(c(1+|u|))\bigr).
\]
We now split this central range into two subranges.

If $|u|\le c^{-1/4}$, then $c|u|^3=O(c^{1/4})$, so the exponential may be linearized and
\[
p_{k,\lambda}
=
c\,\varphi(u)\Bigl(1+O\!\bigl(c(1+|u|^3)\bigr)\Bigr)
=
c\,\varphi(u)+O\!\left(c^2(1+|u|)^{-3}\right),
\]
because $\sup_{u\in\R}\varphi(u)(1+|u|)^6<\infty$.

If instead $c^{-1/4}<|u|\le c^{-1/3}$, then $u^2\ge c^{-1/2}$ and $c|u|^3\le c^{1/4}$. For all sufficiently
small $c$ this gives
\[
p_{k,\lambda}\le Cc\,e^{-u^2/2+c^{1/4}}\le Cc\,e^{-u^2/4}\le Cc\,e^{-c^{-1/2}/4},
\]
while also
\[
c\,\varphi(u)\le Cc\,e^{-u^2/2}\le Cc\,e^{-c^{-1/2}/2}.
\]
Both quantities are super-exponentially small in $c^{-1/2}$, hence in particular they are
$O\!\left(c^2(1+|u|)^{-3}\right)$.

Finally, if $|u|>c^{-1/3}$, then $|k-\lambda|=|u|/c\ge c^{-4/3}=\lambda^{2/3}$. If $k\ge \lambda+\lambda^{2/3}$,
then $p_{k,\lambda}\le \Pbb(\Poi(\lambda)\ge k)$, and the standard one-sided Chernoff bound gives
\[
p_{k,\lambda}\le e^{-c_1\lambda^{1/3}}=e^{-c_1c^{-2/3}}
\]
for some absolute constant $c_1>0$. If $0\le k\le \lambda-\lambda^{2/3}$, then similarly
\[
p_{k,\lambda}\le \Pbb(\Poi(\lambda)\le k)\le e^{-c_1c^{-2/3}}.
\]
Moreover,
\[
c\,\varphi(u)\le c\,e^{-u^2/2}\le c\,e^{-c^{-2/3}/2}.
\]
Again both bounds are $O\!\left(c^2(1+|u|)^{-3}\right)$. This proves
\eqref{eq:poisson-local-weighted}.

Since $|\varphi'(x)|\le C(1+|x|)^{-3}$ for all $x$, integration over the cell $I_{k,c}$ yields
\[
\int_{I_{k,c}}\varphi(x)\,dx
=
c\varphi(z_{k,c})+O\!\left(c^2(1+|z_{k,c}|)^{-3}\right),
\]
and therefore
\begin{equation}
\left|\Pbb(Z_c=z_{k,c})-\int_{I_{k,c}}\varphi(x)\,dx\right|
\le \frac{C_1c^2}{(1+|z_{k,c}|)^3}.
\label{eq:poisson-cell-bound}
\end{equation}

For the deficiency from Poisson to Gaussian, let $K_c$ be the uniform-jitter kernel
\[
K_c(x,\cdot):=\mathrm{Unif}[x-c/2,x+c/2].
\]
The law $K_c\mathcal L(Z_c)$ has density
\[
f_c(x)=\frac{1}{c}\Pbb(Z_c=z_{k,c})\qquad \text{for }x\in I_{k,c}.
\]
Let $\bar\varphi_{k,c}:=c^{-1}\int_{I_{k,c}}\varphi(u)\,du$. Then
\[
\int_{I_{k,c}}|f_c(x)-\varphi(x)|\,dx
\le \left|\Pbb(Z_c=z_{k,c})-\int_{I_{k,c}}\varphi(u)\,du\right|
+\int_{I_{k,c}}|\bar\varphi_{k,c}-\varphi(x)|\,dx.
\]
The first term is bounded by \eqref{eq:poisson-cell-bound}, and the second is at most
$Cc^2(1+|z_{k,c}|)^{-3}$ by the derivative bound on $\varphi$. Hence
\[
\int_{I_{k,c}}|f_c(x)-\varphi(x)|\,dx\le \frac{C_2c^2}{(1+|z_{k,c}|)^3}.
\]
Summing over $k$ and using that $z_{k,c}$ runs over a lattice of mesh $c$,
\[
\TV\bigl(K_c\mathcal L(Z_c),N(0,1)\bigr)
\le C_2c^2\sum_{k\in\mathbb Z}(1+|z_{k,c}|)^{-3}
\le C_3c.
\]
By translation invariance of both the jitter kernel and the Gaussian family,
\[
\TV\bigl(K_c\mathcal L(Z_c+c),N(c,1)\bigr)
=\TV\bigl(K_c\mathcal L(Z_c),N(0,1)\bigr)
\le C_3c.
\]

For the reverse deficiency, let $L_c$ be the rounding kernel that maps $x\in I_{k,c}$ to $z_{k,c}$. Then
$L_cN(0,1)$ is the lattice law assigning mass $\int_{I_{k,c}}\varphi(x)\,dx$ to $z_{k,c}$, so by
\eqref{eq:poisson-cell-bound},
\[
\TV\bigl(L_cN(0,1),\mathcal L(Z_c)\bigr)
\le \frac12\sum_{k\in\mathbb Z}\frac{C_1c^2}{(1+|z_{k,c}|)^3}
\le C_4c.
\]
Again translation invariance yields
\[
\TV\bigl(L_cN(c,1),\mathcal L(Z_c+c)\bigr)\le C_4c.
\]
Taking $C:=\max\{C_3,C_4\}$ proves the theorem.
\end{proof}

\begin{corollary}[Boundary \privacycurve\ convergence]
\label{cor:boundary-BE-privacy}
For every fixed $\varepsilon\ge 0$ there exists $C_\varepsilon<\infty$ such that, for all sufficiently small
$c>0$,
\[
\bigl|\delta_{\mathcal E_c^{\mathrm{Poi}}}(\varepsilon)-\delta_{\mathcal E_c^{\mathrm G}}(\varepsilon)\bigr|
\le C_\varepsilon c,
\]
where $\delta_{\mathcal E}(\varepsilon)$ denotes the privacy curve of the binary experiment $\mathcal E$.
For each fixed $\pi\in(0,1)$, the same conclusion holds for the centered-and-normalized Skellam-shift
experiment of Part~II, with a constant $C_{\varepsilon,\pi}$.
\end{corollary}

\begin{proof}
For the Poisson experiment, combine Theorem~\ref{thm:boundary-BE-poisson} with the monotonicity of privacy curves
under Markov kernels and Lemma~\ref{lem:privacy-stability-main}. Explicitly, let $\eta_c:=Cc$ be the two-sided
Le Cam error from the theorem. Applying $K_c$ to the Poisson experiment and then using \privacycurve\ stability on
the common real line gives
\[
\delta_{\mathcal E_c^{\mathrm{Poi}}}(\varepsilon)
\ge
\delta_{\mathcal E_c^{\mathrm G}}(\varepsilon)-(1+e^\varepsilon)\eta_c.
\]
Applying $L_c$ in the reverse direction gives the opposite inequality. Hence
\[
\bigl|\delta_{\mathcal E_c^{\mathrm{Poi}}}(\varepsilon)-\delta_{\mathcal E_c^{\mathrm G}}(\varepsilon)\bigr|
\le (1+e^\varepsilon)Cc.
\]
This proves the first statement.

For the Skellam case, fix $\pi\in(0,1)$, let
\[
U_c\sim \Poi((1-\pi)c^{-2}),
\qquad
V_c\sim \Poi(\pi c^{-2})
\]
be independent, and set
\[
\widehat S_{c,\pi}:=c\bigl(U_c-V_c-(1-2\pi)c^{-2}\bigr).
\]
Then the centered-and-normalized Skellam-shift experiment from Part~II is
\[
\bigl(\mathcal L(\widehat S_{c,\pi}),\ \mathcal L(\widehat S_{c,\pi}+c)\bigr).
\]
Apply Theorem~\ref{thm:boundary-BE-poisson} to the first Poisson coordinate with parameter
$c/\sqrt{1-\pi}$, and to the second coordinate (which is common under both hypotheses) with parameter
$c/\sqrt{\pi}$. After rescaling those one-dimensional approximations and tensorizing them via
Lemma~\ref{lem:tv-tensorization-main}, we obtain an $O(c)$ comparison between
\[
\bigl(c(U_c-(1-\pi)c^{-2}),\ c(V_c-\pi c^{-2})\bigr)
\]
and the Gaussian product law
\[
N(0,1-\pi)\otimes N(0,\pi),
\]
with the neighboring alternative shifting only the first coordinate by $c$. Pushing forward under the subtraction
map $(x,y)\mapsto x-y$ therefore yields an $O(c)$ comparison between the centered Skellam-shift experiment and
$(N(0,1),N(c,1))$. The same privacy-curve argument as above gives the claimed bound, with a constant depending on
$\pi$.
\end{proof}

\section{Discussion}

\begin{tcolorbox}[keybox,title=What the trilogy now covers]
Parts~I--III together give a finite-alphabet universality theory for neighboring shuffle experiments under
convergent macroscopic scalings. Part~I treats the fixed full-support regime and proves Gaussian/GDP limits with
Berry--Esseen bounds. Part~II treats the critical one-dominant frontier and proves Poisson, Skellam, and
multivariate compound-Poisson limits, together with the super-critical collapse of privacy. The present paper
identifies the general finite-dominant L\'evy--Khintchine layer: a Gaussian factor on the dominant tangent space, a
compound-Poisson jump field on the quotient space, projected Le Cam convergence, and \privacycurve\ convergence for
the full hybrid experiment in the interior, in the weak boundary regime, and in the regular strong-boundary cases.
It also isolates the obstruction showing that a bounded minority block can defeat any fully general projected
strong-boundary \privacycurve\ theorem.
\end{tcolorbox}

Four natural directions remain open. First, the present neighboring calibration produces no nonzero Gaussian shift
in the limit. A genuine hybrid shift theorem
\[
(G,J)\quad\text{versus}\quad (G+h,J+\Delta),\qquad h\neq 0,
\]
would require a simultaneous $n^{-1/2}$ perturbation of the dominant block and would amount to the next L\'evy--
Khintchine layer beyond the neighboring model treated here. Second, the strong-boundary obstruction of
Proposition~\ref{prop:strong-boundary-obstruction} suggests a refined finite-minority theory: when the minority
group has bounded size and retains nontrivial dominant internal geometry, what is the correct limit experiment and
what replaces the projected jump \privacycurve? Third, Part~I treats unbundled multi-message shuffling in the
Gaussian regime, but a critical non-Gaussian theory for unbundled mechanisms remains open. Fourth, the finite-
alphabet assumption in the present paper could likely be relaxed to countable alphabets under summable rare
intensity
\[
\Lambda_b:=\sum_{y\notin D_b}\alpha_b(y)<\infty,
\]
but a full countable-alphabet theory with growing dominant dimension lies beyond the present trilogy. A separate
sharpness question remains open in the overlap regime $\Delta=0$ within the regimes covered by
Theorem~\ref{thm:general-k-dominant-LK}\textup{(iii)}: the present argument gives privacy collapse, but it is not
yet clear whether the optimal rate there is always $n^{-1/2}$ or can be faster.

\appendix
\section{Technical lemmas}
\label{app:technical-lemmas}

All spaces in this appendix are finite-dimensional Euclidean or finite spaces, hence standard Borel. In particular,
all conditional laws and disintegrations used below exist.

\begin{lemma}[Quotient geometry of the dominant block]
\label{lem:new-geometry}
Let $\mathsf H$ be the hypergraph with vertex set $D_0\cup D_1$ and hyperedges $D_0,D_1$, and let $\mathcal C$
denote its connected components. Then
\[
M^\perp\cap \mathbb R^{D_0\cup D_1}
=
\{x\in \mathbb R^{D_0\cup D_1}:\ x \text{ is constant on each } C\in\mathcal C\}.
\]
Equivalently,
\[
\Pi_J e_y=\Pi_J e_{y'}
\quad\Longleftrightarrow\quad
y,y' \text{ belong to the same component } C\in\mathcal C.
\]
In particular, for each $b\in\{0,1\}$ and each $y\in D_b$,
\[
\Pi_J e_y=\Pi_J\mu_b=:m_b.
\]
Hence all dominant outputs in $D_b$ collapse to a single quotient atom $m_b\in M^\perp$.

Moreover:
\begin{enumerate}
\item[(a)] if $D_0\cap D_1=\varnothing$, then $m_0\neq m_1$;
\item[(b)] if $D_0\cap D_1\neq\varnothing$, then $m_0=m_1$, and therefore
\[
\Delta=\Pi_J(\mu_1-\mu_0)=0.
\]
\end{enumerate}
Finally, for every $y\notin D_b$,
\[
j_{b,y}=\Pi_J(e_y-\mu_b)=\Pi_J e_y-m_b.
\]
\end{lemma}

\begin{proof}
Let $U:=\mathbb R^{D_0\cup D_1}$. Since
\[
M_b=\{x\in U:\ \supp(x)\subseteq D_b,\ \sum_{y\in D_b}x(y)=0\},
\]
a vector $x\in U$ belongs to $M^\perp\cap U$ if and only if it is orthogonal to every zero-sum vector on $D_0$ and
every zero-sum vector on $D_1$. Equivalently, $x$ is constant on $D_0$ and constant on $D_1$. If $D_0\cap D_1=
\varnothing$, the connected components of $\mathsf H$ are $D_0$ and $D_1$, so this is exactly the statement that $x$
is constant on each component. If $D_0\cap D_1\neq\varnothing$, the two constants agree on the overlap, hence $x$ is
constant on the single component $D_0\cup D_1$. This proves the first display.

Fix $b\in\{0,1\}$ and $y\in D_b$. Since $e_y-\mu_b\in M_b\subseteq M$, we have
\[
\Pi_J(e_y-\mu_b)=0,
\]
which is equivalent to
\[
\Pi_J e_y=\Pi_J\mu_b=:m_b.
\]
Thus all dominant outputs in $D_b$ collapse to the same quotient atom.

If $D_0\cap D_1\neq\varnothing$, choose $s\in D_0\cap D_1$. Then
\[
m_0=\Pi_J e_s=m_1,
\]
so
\[
\Delta=\Pi_J(\mu_1-\mu_0)=m_1-m_0=0.
\]
In this case $\mathsf H$ has a single component, so the equivalence for $\Pi_J e_y$ follows immediately.

Assume now that $D_0\cap D_1=\varnothing$ and suppose for contradiction that $m_0=m_1$. Choose any
$y_0\in D_0$ and $y_1\in D_1$. Then
\[
\Pi_J(e_{y_0}-e_{y_1})=0,
\]
so $e_{y_0}-e_{y_1}\in M=M_0+M_1$. Hence there exist $x_0\in M_0$ and $x_1\in M_1$ with
\[
e_{y_0}-e_{y_1}=x_0+x_1.
\]
Because $D_0$ and $D_1$ are disjoint, the supports of $x_0$ and $x_1$ are disjoint. Restricting to $D_0$ gives
$x_0=e_{y_0}$, while restricting to $D_1$ gives $x_1=-e_{y_1}$. But vectors in $M_b$ have coordinate sum zero on
$D_b$, whereas $e_{y_0}$ and $-e_{y_1}$ have sums $1$ and $-1$, respectively. This contradiction proves $m_0\neq
m_1$, establishing part~(a). Since the components are now exactly $D_0$ and $D_1$, the displayed equivalence for
$\Pi_J e_y$ follows in the disjoint case as well.

Finally, for $y\notin D_b$,
\[
j_{b,y}=\Pi_J(e_y-\mu_b)=\Pi_J e_y-\Pi_J\mu_b=\Pi_J e_y-m_b.
\]
\end{proof}

\begin{lemma}[One-user characteristic expansion]
\label{lem:new-one-user-cf}
Fix $b\in\{0,1\}$ and define
\[
\phi_{b,n}(u,v)
:=
\mathbb E\exp\!\left(
 i\bigl\langle u,n^{-1/2}\Pi_G X^{(b)}_{1,n}\bigr\rangle
 +i\bigl\langle v,\Pi_J X^{(b)}_{1,n}\bigr\rangle
\right),
\qquad (u,v)\in M\times M^\perp.
\]
Then, uniformly on compact subsets of $M\times M^\perp$,
\[
\phi_{b,n}(u,v)
=
1
-\frac{1}{2n}\langle u,\Gamma_b u\rangle
+\frac{1}{n}\sum_{y\notin D_b}\alpha_b(y)\bigl(e^{i\langle v,j_{b,y}\rangle}-1\bigr)
+o(n^{-1}),
\]
and hence
\[
\log\phi_{b,n}(u,v)
=
-\frac{1}{2n}\langle u,\Gamma_b u\rangle
+\frac{1}{n}\sum_{y\notin D_b}\alpha_b(y)\bigl(e^{i\langle v,j_{b,y}\rangle}-1\bigr)
+o(n^{-1}).
\]
\end{lemma}

\begin{proof}
Fix a compact set $K\subset M\times M^\perp$. All $O(\cdot)$ and $o(\cdot)$ terms below are uniform in $(u,v)\in K$.
Write
\[
X_n:=X^{(b)}_{1,n}=e_{Y_n}-\mu_b,
\qquad Y_n\sim W_b^{(n)}.
\]
Split the expectation into dominant and rare outputs.

If $y\in D_b$, then $e_y-\mu_b\in M_b\subseteq M$, so $\Pi_G(e_y-\mu_b)=e_y-\mu_b$ and
$\Pi_J(e_y-\mu_b)=0$. Therefore the dominant contribution equals
\[
D_n(u):=\sum_{y\in D_b} W_b^{(n)}(y)\exp\!\bigl(i n^{-1/2}\langle u,e_y-\mu_b\rangle\bigr).
\]
By assumption,
\[
W_b^{(n)}(y)=p_b(y)+r_{b,n}(y),\qquad r_{b,n}(y)=O(n^{-1}),\qquad y\in D_b,
\]
and because the rare mass is $O(n^{-1})$,
\[
\sum_{y\in D_b}p_b(y)=1,
\qquad
\sum_{y\in D_b}W_b^{(n)}(y)=1+O(n^{-1}).
\]
Using the Taylor expansion
\[
e^{it}=1+it-\frac{t^2}{2}+O(|t|^3)
\]
with $t=n^{-1/2}\langle u,e_y-\mu_b\rangle$, we obtain
\begin{align*}
D_n(u)
&=
\sum_{y\in D_b} W_b^{(n)}(y)
\left(1+i n^{-1/2}\langle u,e_y-\mu_b\rangle-
\frac{1}{2n}\langle u,e_y-\mu_b\rangle^2+O(n^{-3/2})\right) \\
&=
\sum_{y\in D_b}W_b^{(n)}(y)
+i n^{-1/2}\Bigl\langle u,\sum_{y\in D_b}W_b^{(n)}(y)(e_y-\mu_b)\Bigr\rangle \\
&\qquad-
\frac{1}{2n}\sum_{y\in D_b}W_b^{(n)}(y)\langle u,e_y-\mu_b\rangle^2+O(n^{-3/2}).
\end{align*}
Now
\[
\sum_{y\in D_b}W_b^{(n)}(y)(e_y-\mu_b)
=
\sum_{y\in D_b}(W_b^{(n)}(y)-p_b(y))(e_y-\mu_b)=O(n^{-1}),
\]
so the linear term is $O(n^{-3/2})$. Also,
\[
\sum_{y\in D_b}W_b^{(n)}(y)\langle u,e_y-\mu_b\rangle^2
=
\sum_{y\in D_b}p_b(y)\langle u,e_y-\mu_b\rangle^2+O(n^{-1})
=\langle u,\Gamma_bu\rangle+O(n^{-1}),
\]
whence
\[
D_n(u)=1-\sum_{y\notin D_b}W_b^{(n)}(y)-\frac{1}{2n}\langle u,\Gamma_bu\rangle+o(n^{-1}).
\]
Since $nW_b^{(n)}(y)\to\alpha_b(y)$ for $y\notin D_b$ and the alphabet is finite,
\[
\sum_{y\notin D_b}W_b^{(n)}(y)=\frac{1}{n}\sum_{y\notin D_b}\alpha_b(y)+o(n^{-1}).
\]
Therefore
\[
D_n(u)=1-\frac{1}{n}\sum_{y\notin D_b}\alpha_b(y)-\frac{1}{2n}\langle u,\Gamma_bu\rangle+o(n^{-1}).
\]

For $y\notin D_b$, the contribution equals
\[
R_{n,y}(u,v):=W_b^{(n)}(y)
\exp\!\left(i n^{-1/2}\langle u,\Pi_G(e_y-\mu_b)\rangle+i\langle v,j_{b,y}\rangle\right).
\]
Because $\Pi_G(e_y-\mu_b)$ is fixed and $(u,v)$ stays in a compact set,
\[
\exp\!\left(i n^{-1/2}\langle u,\Pi_G(e_y-\mu_b)\rangle+i\langle v,j_{b,y}\rangle\right)
=
 e^{i\langle v,j_{b,y}\rangle}\bigl(1+O(n^{-1/2})\bigr).
\]
Multiplying by $W_b^{(n)}(y)=\alpha_b(y)n^{-1}+o(n^{-1})$ yields
\[
R_{n,y}(u,v)=\frac{1}{n}\alpha_b(y)e^{i\langle v,j_{b,y}\rangle}+o(n^{-1}).
\]
Summing over all rare outputs gives
\[
\sum_{y\notin D_b}R_{n,y}(u,v)
=
\frac{1}{n}\sum_{y\notin D_b}\alpha_b(y)e^{i\langle v,j_{b,y}\rangle}+o(n^{-1}).
\]
Combining the dominant and rare parts proves the first display.

Since $\phi_{b,n}(u,v)=1+O(n^{-1})$ uniformly on $K$, the logarithm satisfies
\[
\log\phi_{b,n}(u,v)=\phi_{b,n}(u,v)-1+o(n^{-1}),
\]
which yields the second display.
\end{proof}

\begin{lemma}[Projected quotient array]
\label{lem:new-projected-array}
Let
\[
\mathcal Y_J:=\Pi_J\bigl(\{e_y:y\in\mathcal Y\}\bigr)\subset M^\perp,
\]
and let $\widetilde W_b^{(n)}$ be the pushforward of $W_b^{(n)}$ under $y\mapsto \Pi_J e_y$.
Then, for each $b\in\{0,1\}$,
\[
\widetilde W_b^{(n)}(m_b)=1-O(n^{-1}),
\]
and for every $w\in\mathcal Y_J\setminus\{m_b\}$,
\[
n\,\widetilde W_b^{(n)}(w)\longrightarrow \widetilde\alpha_b(w)
:=
\sum_{\substack{y\notin D_b\\ \Pi_J e_y=w}}\alpha_b(y).
\]
Consequently:
\begin{enumerate}
\item[(a)] if $m_0\neq m_1$, then $\widetilde W^{(n)}$ is exactly in the sparse-error critical regime of
Part~II, Definition~5.1 on the finite alphabet $\mathcal Y_J$;
\item[(b)] if $m_0=m_1$, then $\widetilde W^{(n)}$ is in the coincident-dominant variant of that regime,
with a common dominant atom and zero deterministic shift.
\end{enumerate}
\end{lemma}

\begin{proof}
Fix $b\in\{0,1\}$. By Lemma~\ref{lem:new-geometry}, all points of $D_b$ map to the same quotient atom
$m_b=\Pi_Je_y=\Pi_J\mu_b$. Therefore
\[
\widetilde W_b^{(n)}(m_b)
=
\sum_{y\in D_b}W_b^{(n)}(y)+\sum_{\substack{y\notin D_b\\ \Pi_Je_y=m_b}}W_b^{(n)}(y).
\]
The first sum equals $1-O(n^{-1})$ because the total rare mass is $O(n^{-1})$, and the second sum is itself
$O(n^{-1})$, so
\[
\widetilde W_b^{(n)}(m_b)=1-O(n^{-1}).
\]

Now fix $w\neq m_b$. If $y\in D_b$, then $\Pi_Je_y=m_b$, so no dominant output contributes to $w$. Hence
\[
\widetilde W_b^{(n)}(w)=\sum_{\substack{y\notin D_b\\ \Pi_Je_y=w}}W_b^{(n)}(y).
\]
Multiplying by $n$ and using the pointwise limits $nW_b^{(n)}(y)\to \alpha_b(y)$ yields
\[
n\,\widetilde W_b^{(n)}(w)
\to
\sum_{\substack{y\notin D_b\\ \Pi_Je_y=w}}\alpha_b(y)=\widetilde\alpha_b(w).
\]
This proves the displayed limits.

If $m_0\neq m_1$, then the projected array has two distinct dominant quotient atoms and rare probabilities of
order $1/n$ away from them, exactly as in Part~II, Definition~5.1. If $m_0=m_1$, then the projected array has a
single common dominant quotient atom and the same rare-$1/n$ behavior, but the deterministic neighboring shift
vanishes because $\Delta=m_1-m_0=0$. This proves parts~(a) and~(b).
\end{proof}

\begin{lemma}[Measurability of dominant totals in the disjoint quotient case]
\label{lem:new-measurable-totals}
Assume $m_0\neq m_1$. Then there exist measurable maps
\[
\tau_{b,n}:M^\perp\to \mathbb N,\qquad b\in\{0,1\},
\]
such that
\[
N_{n,k_n}(D_b)=\tau_{b,n}(Z_n)\qquad\text{a.s. under both hypotheses.}
\]
Equivalently, the dominant totals on $D_0$ and $D_1$ are $\sigma(Z_n)$-measurable.
\end{lemma}

\begin{proof}
Let
\[
\Omega_n:=\Bigl\{h\in\mathbb Z^{\mathcal Y}:\sum_{y\in\mathcal Y}h(y)=n\Bigr\}
\]
be the finite set of histograms of size $n$. For $h\in\Omega_n$, define
\[
\widehat h:=h-m_{0,n}\mu_0-m_{1,n}\mu_1\in\R^{\mathcal Y},
\qquad z(h):=\Pi_J\widehat h\in M^\perp.
\]
Suppose $h,h'\in\Omega_n$ satisfy $z(h)=z(h')$. Then
\[
\Pi_J(\widehat h-\widehat h')=0,
\]
so $\widehat h-\widehat h'\in M$. Since $D_0$ and $D_1$ are disjoint in the present case (by Lemma~\ref{lem:new-geometry},
$m_0\neq m_1$ implies $D_0\cap D_1=\varnothing$), every vector $x\in M=M_0+M_1$ decomposes as $x=x_0+x_1$ with
$x_b\in M_b$, $\supp(x_b)\subseteq D_b$, and
\[
\sum_{y\in D_b}x_b(y)=0.
\]
Therefore
\[
\sum_{y\in D_b}(\widehat h(y)-\widehat h'(y))=0,
\qquad b\in\{0,1\}.
\]
Because the centering vectors $m_{0,n}\mu_0+m_{1,n}\mu_1$ have the same totals on $D_b$ for both $h$ and $h'$, this
is equivalent to
\[
h(D_b)=h'(D_b),\qquad b\in\{0,1\}.
\]
Hence the quantity $h(D_b)$ depends only on the fiber value $z(h)$. Since $\Omega_n$ is finite, we may define
$\tau_{b,n}(z)$ to be this common value on the fiber $\{h\in\Omega_n:z(h)=z\}$ and set $\tau_{b,n}(z)=0$ on points
outside the support of $Z_n$. The resulting map is measurable because the support is finite. Evaluating at the
random histogram $N_{n,k_n}$ proves the claim.
\end{proof}

\begin{lemma}[Multinomial edge-shift bound]
\label{lem:new-edge-shift}
Let $D$ be finite, let $\vartheta\in\Delta(D)$ satisfy
\[
\vartheta_{\min}:=\min_{y\in D}\vartheta(y)>0,
\]
and let $S_m\sim \mathrm{Mult}(m,\vartheta)$.
Then there exists $C=C(\vartheta)<\infty$ such that for every $a,b\in D$,
\[
\mathrm{TV}\bigl(\mathcal L(S_m+e_a),\mathcal L(S_m+e_b)\bigr)
\le \frac{C}{\sqrt{m+1}}.
\]
More generally, for deterministic $a_1,\dots,a_r,b_1,\dots,b_r\in D$,
\[
\mathrm{TV}\!\left(
\mathcal L\!\left(S_m+\sum_{j=1}^r e_{a_j}\right),
\mathcal L\!\left(S_m+\sum_{j=1}^r e_{b_j}\right)
\right)
\le \frac{Cr}{\sqrt{m+1}}.
\]
\end{lemma}

\begin{proof}
We first treat a single shift. Fix $a,b\in D$, $a\neq b$. Condition on the vector of counts outside $\{a,b\}$,
namely
\[
V:=(S_m(y))_{y\in D\setminus\{a,b\}}.
\]
Let $R:=\sum_{y\in D\setminus\{a,b\}}V(y)$. Given $V$, the remaining pair $(S_m(a),S_m(b))$ is distributed as
\[
(B,m-R-B),\qquad B\sim \Bin(m-R,p_{ab}),
\qquad p_{ab}:=\frac{\vartheta(a)}{\vartheta(a)+\vartheta(b)}.
\]
Therefore
\[
\TV\bigl(\mathcal L(S_m+e_a\mid V),\mathcal L(S_m+e_b\mid V)\bigr)
=
\TV\bigl(\mathcal L(B+1),\mathcal L(B)\bigr).
\]
For any integer-valued random variable $B$ with unimodal pmf,
\[
\TV\bigl(\mathcal L(B+1),\mathcal L(B)\bigr)=\max_k \Pbb(B=k).
\]
Since binomial laws are unimodal, we get
\[
\TV\bigl(\mathcal L(B+1),\mathcal L(B)\bigr)
\le \frac{c_{ab}}{\sqrt{m-R+1}},
\qquad
c_{ab}:=\sqrt{\frac{2}{p_{ab}(1-p_{ab})}}.
\]
Taking expectation over $V$ yields
\[
\TV\bigl(\mathcal L(S_m+e_a),\mathcal L(S_m+e_b)\bigr)
\le c_{ab}\,\E\Bigl[(m-R+1)^{-1/2}\Bigr].
\]
Now $m-R=S_m(a)+S_m(b)\sim \Bin(m,\vartheta(a)+\vartheta(b))$. Write
\[
\rho_{ab}:=\vartheta(a)+\vartheta(b)\ge 2\vartheta_{\min}>0.
\]
On the event $\{m-R\ge \rho_{ab}m/2\}$,
\[
(m-R+1)^{-1/2}\le \sqrt{\frac{2}{\rho_{ab}m}}.
\]
On the complementary event, the quantity is bounded by $1$. A Chernoff bound gives
\[
\Pbb\Bigl(m-R<\frac{\rho_{ab}m}{2}\Bigr)\le e^{-c'm}
\]
for some $c'=c'(\vartheta)>0$. Consequently,
\[
\E\Bigl[(m-R+1)^{-1/2}\Bigr]
\le \sqrt{\frac{2}{\rho_{ab}m}}+e^{-c'm}
\le \frac{C_{ab}}{\sqrt{m+1}}.
\]
Since $D$ is finite, taking the maximum over all pairs $(a,b)$ gives a constant $C=C(\vartheta)$ such that
\[
\TV\bigl(\mathcal L(S_m+e_a),\mathcal L(S_m+e_b)\bigr)
\le \frac{C}{\sqrt{m+1}}
\]
for all $a,b\in D$.

For the general $r$-shift bound, define intermediate sums
\[
T^{(0)}:=S_m+\sum_{j=1}^r e_{a_j},
\qquad
T^{(r)}:=S_m+\sum_{j=1}^r e_{b_j},
\]
and
\[
T^{(s)}:=S_m+\sum_{j=1}^{s}e_{b_j}+\sum_{j=s+1}^{r}e_{a_j},
\qquad 0\le s\le r.
\]
For each $0\le s\le r-1$, write
\[
h_s:=\sum_{j=1}^{s}e_{b_j}+\sum_{j=s+2}^{r}e_{a_j}.
\]
Then
\[
T^{(s)}=(S_m+e_{a_{s+1}})+h_s,
\qquad
T^{(s+1)}=(S_m+e_{b_{s+1}})+h_s.
\]
Since total variation is invariant under a common deterministic translation, the one-step bound implies
\[
\TV\bigl(\mathcal L(T^{(s)}),\mathcal L(T^{(s+1)})\bigr)
\le \frac{C}{\sqrt{m+1}}.
\]
Therefore the triangle inequality yields
\[
\TV\bigl(\mathcal L(T^{(0)}),\mathcal L(T^{(r)})\bigr)
\le \sum_{s=0}^{r-1}\TV\bigl(\mathcal L(T^{(s)}),\mathcal L(T^{(s+1)})\bigr)
\le \frac{Cr}{\sqrt{m+1}}.
\]
\end{proof}

\begin{lemma}[Conditional multinomial smoothing in the disjoint case]
\label{lem:new-conditional-smoothing}
Assume $\pi\in(0,1)$ and $D_0\cap D_1=\varnothing$.
Define the normalized dominant laws
\[
\vartheta_{b,n}(y):=\frac{W_b^{(n)}(y)}{W_b^{(n)}(D_b)},
\qquad y\in D_b.
\]
For $z$ in the support of $Z_n$, let
\[
T_{b,n}(z):=N_{n,k_n}(D_b),
\]
and let $X_{b,n}^\circ\sim \mathrm{Mult}(T_{b,n}(z),\vartheta_{b,n})$ be independent.
Let $R_{n,z}$ be the law of
\[
\Psi_{n,z}(X_{0,n}^\circ,X_{1,n}^\circ)
:=
n^{-1/2}\Pi_G\!\left(
\sum_{y\in D_0}(X_{0,n}^\circ(y)-T_{0,n}(z)\vartheta_{0,n}(y))e_y
+
\sum_{y\in D_1}(X_{1,n}^\circ(y)-T_{1,n}(z)\vartheta_{1,n}(y))e_y
\right).
\]
Then there exists $C_{\mathrm{int}}<\infty$ such that
\[
\int \mathrm{TV}\bigl(P_{n,z}^{G\mid Z},R_{n,z}\bigr)\,P_n^J(dz)
+
\int \mathrm{TV}\bigl(Q_{n,z}^{G\mid Z},R_{n,z}\bigr)\,Q_n^J(dz)
\le \frac{C_{\mathrm{int}}}{\sqrt n}.
\]
\end{lemma}

\begin{proof}
We prove the bound under an arbitrary one of the two hypotheses; the same argument applies to both, and the two
bounds can then be added. Write again $(m_{0,n}^\star,m_{1,n}^\star)$ for the actual group sizes under the chosen
hypothesis.

Define $L_{b,n}$ to be the number of users from group $b$ whose outputs leave $D_b$, and define $A_{b,n}$ to be the
number of users from the opposite group landing inside $D_b$. Because $D_0$ and $D_1$ are disjoint, every cross
message is rare, so $A_{b,n}\le L_{1-b,n}$. Refine the conditioning exactly as in Section~\ref{sec:proof-iii}: reveal
which users are rare, whether each rare user lands in the opposite dominant block or elsewhere, and reveal all rare
outputs outside $D_0\cup D_1$. Conditional on this refined sigma-field $\mathcal F_n$, the two dominant blocks are
independent, and for each $b\in\{0,1\}$ the dominant count vector on $D_b$ has the form
\[
Y_{b,n}:=S_{b,n}+\sum_{r=1}^{A_{b,n}}\eta_{b,r},
\]
where
\[
S_{b,n}\sim \Mult(m_{b,n}^\star-L_{b,n},\vartheta_{b,n})
\]
is the native dominant contribution, while $\eta_{b,1},\dots,\eta_{b,A_{b,n}}$ are i.i.d. categorical vectors in
$\{e_y:y\in D_b\}$ with some cross split law $\rho_{b,n}$ determined by the opposite group. The surrogate block is
\[
Y_{b,n}^\circ:=S_{b,n}+\sum_{r=1}^{A_{b,n}}\eta_{b,r}^\circ,
\]
where the $\eta_{b,r}^\circ$ are i.i.d. $\Cat(\vartheta_{b,n})$, independent of everything else. Conditional on
$Z_n=z$, the law of the pair $(Y_{0,n}^\circ,Y_{1,n}^\circ)$ is exactly the pair of independent multinomials used to
define $R_{n,z}$, because
\[
T_{b,n}(z)=(m_{b,n}^\star-L_{b,n})+A_{b,n}.
\]

Fix $b$ and condition on $\mathcal F_n$. By the same telescoping argument as in Section~\ref{sec:proof-iii}, define
interpolating sums replacing the cross messages one by one. Lemma~\ref{lem:new-edge-shift} yields
\[
\TV\bigl(\mathcal L(Y_{b,n}\mid \mathcal F_n),\mathcal L(Y_{b,n}^\circ\mid \mathcal F_n)\bigr)
\le \frac{C_bA_{b,n}}{\sqrt{m_{b,n}^\star-L_{b,n}+1}}.
\]
Since the two blocks are conditionally independent and the map from dominant count vectors to the Gaussian
coordinate is measurable, Lemma~\ref{lem:tv-tensorization-main} implies
\begin{align*}
\TV\bigl(\mathcal L(G_n\mid Z_n,\mathcal F_n),R_{n,Z_n}\bigr)
&\le \frac{C_0A_{0,n}}{\sqrt{m_{0,n}^\star-L_{0,n}+1}}
+\frac{C_1A_{1,n}}{\sqrt{m_{1,n}^\star-L_{1,n}+1}}.
\end{align*}
Averaging over $\mathcal F_n$ and then over $Z_n$ gives the same bound for the regular conditional laws
$P_{n,z}^{G\mid Z}$ or $Q_{n,z}^{G\mid Z}$.

It remains to estimate the expectation of the right-hand side. Because $\pi\in(0,1)$ there exists $\kappa>0$ such
that, for all sufficiently large $n$ and under both hypotheses,
\[
m_{0,n}^\star\ge 3\kappa n,
\qquad
m_{1,n}^\star\ge 3\kappa n.
\]
Let
\[
R_n:=L_{0,n}+L_{1,n}.
\]
Since the rare probabilities are $O(n^{-1})$ and the alphabet is finite, $R_n$ is a Poisson-binomial sum with
bounded first and second moments:
\[
\E[R_n]\le C,
\qquad
\E[R_n^2]\le C.
\]
On the event $\{R_n\le \kappa n\}$,
\[
\frac{A_{b,n}}{\sqrt{m_{b,n}^\star-L_{b,n}+1}}
\le \frac{A_{b,n}}{\sqrt{2\kappa n}}.
\]
On the complementary event the same quantity is bounded by $A_{b,n}\le R_n$. Hence
\begin{align*}
\E\Bigl[\frac{A_{b,n}}{\sqrt{m_{b,n}^\star-L_{b,n}+1}}\Bigr]
&\le \frac{\E[A_{b,n}]}{\sqrt{2\kappa n}}+\E\bigl[R_n\1\{R_n>\kappa n\}\bigr] \\
&\le \frac{C}{\sqrt n}+\frac{\E[R_n^2]}{\kappa n}
\le \frac{C'}{\sqrt n}.
\end{align*}
Summing over $b=0,1$ proves the stated $O(n^{-1/2})$ bound.
\end{proof}

\begin{lemma}[Overlap-case one-user replacement bound]
\label{lem:new-overlap-tv}
Assume $\pi\in(0,1)$ and $D_0\cap D_1\neq\varnothing$.
Then there exists $C<\infty$ such that
\[
\mathrm{TV}(P_n,Q_n)\le \frac{C}{\sqrt n}.
\]
Consequently, for every fixed $\varepsilon\ge 0$,
\[
0\le \delta_{Q_n\|P_n}(\varepsilon)\le \frac{C}{\sqrt n}.
\]
\end{lemma}

\begin{proof}
Choose a common dominant symbol $s\in D_0\cap D_1$. Couple the two neighboring experiments by using the same
$n-1$ common users under both hypotheses and only switching the last user from input $0$ to input $1$. Let
$H_n^{\mathrm{com}}$ denote the centered contribution of the common users under the centering
$(n-k_n-1)\mu_0+k_n\mu_1$. Let $A_n\sim W_0^{(n)}$ and $B_n\sim W_1^{(n)}$ be the outputs of the switched user under the
null and alternative, respectively, independent of the common part. Then
\[
S_n^{(P)}=H_n^{\mathrm{com}}+\Xi(A_n),
\qquad
S_n^{(Q)}=H_n^{\mathrm{com}}+\Xi(B_n),
\]
where
\[
\Xi(y):=\bigl(n^{-1/2}\Pi_G(e_y-\mu_0),\ \Pi_J(e_y-\mu_0)\bigr).
\]
Therefore
\[
\TV(P_n,Q_n)
\le
\TV\bigl(\mathcal L(H_n^{\mathrm{com}}+\Xi(A_n)),\mathcal L(H_n^{\mathrm{com}}+\Xi(s))\bigr)
+
\TV\bigl(\mathcal L(H_n^{\mathrm{com}}+\Xi(s)),\mathcal L(H_n^{\mathrm{com}}+\Xi(B_n))\bigr).
\]
We treat the first term; the second is identical with the roles of the two groups interchanged.

Split according to whether the switched user is rare. Since $W_0^{(n)}(\mathcal Y\setminus D_0)=O(n^{-1})$,
\begin{align*}
&\TV\bigl(\mathcal L(H_n^{\mathrm{com}}+\Xi(A_n)),\mathcal L(H_n^{\mathrm{com}}+\Xi(s))\bigr) \\
&\le W_0^{(n)}(\mathcal Y\setminus D_0)
+\TV\bigl(\mathcal L(H_n^{\mathrm{com}}+\Xi(A_n)\mid A_n\in D_0),\mathcal L(H_n^{\mathrm{com}}+\Xi(s))\bigr).
\end{align*}
Conditional on $A_n\in D_0$, the switched user's location inside $D_0$ has law $\vartheta_{0,n}$. Now condition on all
common $1$-users, on all rare outputs of the common $0$-users, and on the total number of common $0$-users whose
outputs remain in $D_0$. Given this conditioning, the common $0$-block is distributed as
\[
S_{0,n}\sim \Mult(m_{0,n}'-L_{0,n},\vartheta_{0,n})
\]
for some $m_{0,n}'\asymp n$, and the rest of the experiment contributes a common additive term. Therefore,
Lemma~\ref{lem:new-edge-shift} gives
\[
\TV\bigl(\mathcal L(H_n^{\mathrm{com}}+\Xi(A_n)\mid \mathcal F_n,A_n\in D_0),\mathcal L(H_n^{\mathrm{com}}+\Xi(s)\mid\mathcal F_n)\bigr)
\le \frac{C_0}{\sqrt{m_{0,n}'-L_{0,n}+1}}.
\]
Averaging over the conditioning variables and using the same rare-count estimate as in
Lemma~\ref{lem:new-conditional-smoothing} yields an $O(n^{-1/2})$ bound. The same reasoning for the second term,
with the common $1$-block and the switched output $B_n$, gives another $O(n^{-1/2})$ bound. Hence
\[
\TV(P_n,Q_n)\le \frac{C}{\sqrt n}.
\]
The \privacycurve\ statement follows from the elementary bound
\[
0\le \delta_{Q_n\|P_n}(\varepsilon)\le \TV(P_n,Q_n).
\]
\end{proof}

\begingroup
\sloppy
\hbadness=10000
\begin{thebibliography}{99}

\bibitem{Shv26I}
A.~Shvets.
\textit{Universal Shuffle Asymptotics: Sharp Privacy Analysis in the Gaussian Regime}.
arXiv:2602.09029, 2026.

\bibitem{Shv26II}
A.~Shvets.
\textit{Universal Shuffle Asymptotics, Part II: Non-Gaussian Limits for Shuffle Privacy---Poisson, Skellam, and Compound-Poisson Regimes}.
arXiv:2603.10073, 2026.

\bibitem{LeC86}
L.~Le~Cam.
\textit{Asymptotic Methods in Statistical Decision Theory}.
Springer, 1986.

\bibitem{vdV98}
A.~W.~van der Vaart.
\textit{Asymptotic Statistics}.
Cambridge University Press, 1998.

\bibitem{Str85}
H.~Strasser.
Scale invariance of statistical experiments.
\textit{Probability and Mathematical Statistics}, 5:1--20, 1985.

\bibitem{JMS85}
A.~Janssen, H.~Milbrodt, and H.~Strasser.
\textit{Infinitely Divisible Statistical Experiments}.
Lecture Notes in Statistics 27, Springer, 1985.

\bibitem{Car02}
A.~V.~Carter.
Deficiency distance between multinomial and multivariate normal experiments.
\textit{Annals of Statistics}, 30(3):708--730, 2002.

\bibitem{BHJ92}
A.~D.~Barbour, L.~Holst, and S.~Janson.
\textit{Poisson Approximation}.
Oxford University Press, 1992.

\bibitem{Roos19}
B.~Roos.
Poisson approximation.
\textit{Probability Surveys}, 16:1--119, 2019.

\bibitem{DKT15}
C.~Daskalakis, G.~Kamath, and C.~Tzamos.
On the structure, covering, and learning of Poisson multinomial distributions.
In \textit{Proceedings of the 56th IEEE Symposium on Foundations of Computer Science (FOCS)}, 2015.

\bibitem{Pet75}
V.~V.~Petrov.
\textit{Sums of Independent Random Variables}.
Springer, 1975.

\bibitem{DRS22}
J.~Dong, A.~Roth, and W.~J.~Su.
Gaussian differential privacy.
\textit{Journal of the Royal Statistical Society: Series B}, 84(1):3--37, 2022.

\bibitem{EFM19}
\'U.~Erlingsson, V.~Feldman, I.~Mironov, A.~Raghunathan, K.~Talwar, and A.~Thakurta.
Amplification by shuffling: From local to central differential privacy via anonymity.
In \textit{SODA}, 2019.

\bibitem{CSU19}
A.~Cheu, A.~Smith, J.~Ullman, D.~Zeber, and M.~Zhilyaev.
Distributed differential privacy via shuffling.
In \textit{EUROCRYPT}, 2019.

\bibitem{BBG19}
B.~Balle, J.~Bell, A.~Gascon, and K.~Nissim.
The privacy blanket of the shuffle model.
In \textit{CRYPTO}, 2019.

\bibitem{FMT21}
V.~Feldman, A.~McMillan, and K.~Talwar.
Hiding among the clones: A simple and nearly optimal analysis of privacy amplification by shuffling.
In \textit{FOCS}, 2021.

\bibitem{FMT23}
V.~Feldman, A.~McMillan, and K.~Talwar.
Stronger privacy amplification by shuffling for R\'enyi and approximate differential privacy.
In \textit{Proceedings of the 2023 ACM-SIAM Symposium on Discrete Algorithms (SODA)}, pages 4966--4981, 2023.

\bibitem{TL26}
S.~Takagi and S.~P.~Liew.
Analysis of shuffling beyond pure local differential privacy.
\textit{arXiv:2601.19154}, 2026.

\end{thebibliography}
\endgroup

\end{document}
